\documentclass[a4paper]{article}
\usepackage[utf8]{inputenc} % standard unicode
\usepackage[italian]{babel} % corretta sillabazione in italiano
\usepackage{geometry} % per impostare margini e layout pagina
\usepackage{amssymb} % per l'ambiente matematico
\usepackage{amsmath} % per l'ambiente matematico
\usepackage{enumitem} % per elenchi puntati
\usepackage{multirow} % per celle che si espandono su più righe
\usepackage{tabularx} % per tabelle con larghezza flessibile
\usepackage{hyperref} % per link e riferimenti ipertestuali
\usepackage{booktabs} % per linee orizzontali tabelle
\usepackage{diagbox} % per celle con testo diagonale
\usepackage{xcolor} % per colorare testo e sfondi
\usepackage{algorithm} % per codice
\usepackage[noend]{algpseudocode}

% per margini
\geometry{a4paper,left=25mm, right=25mm, bottom=25mm, top=30mm}

% per centrare testo nelle tabelleX
\renewcommand\tabularxcolumn[1]{m{#1}}
\newcolumntype{L}{>{\raggedright\arraybackslash}X}
\newcolumntype{C}{>{\centering\arraybackslash}X}
\newcolumntype{R}{>{\raggedleft\arraybackslash}X}

% operatore argmin
\DeclareMathOperator*{\argmin}{argmin}
\DeclareMathOperator*{\argmax}{argmax}

% per elenchi puntati
\setlist[itemize]{label=-, partopsep=0pt, topsep=3pt, itemsep=0pt}
\setlist[enumerate]{partopsep=0pt, topsep=3pt, itemsep=0pt}

% percorso delle immagini da inserire
%\graphicspath{{./}}

% per algoritmi in pseudocodice
\renewcommand{\algorithmicprocedure}{\textbf{Algoritmo:}\;}
\renewcommand{\algorithmicrequire}{\qquad\textbf{Input:} \quad\;\;\,}
\renewcommand{\algorithmicensure}{\qquad\textbf{Output:} \;\;\;}
\algnewcommand\algorithmicin{\textbf{in}}
\algnewcommand\algorithmicto{\textbf{to}}
%\algrenewtext{For}[3] {\algorithmicfor\ #1 \gets #2 \algorithmicto\ #3 \algorithmicdo}
\algrenewtext{For}[2]{\algorithmicfor\ #1 \algorithmicto\ #2 \algorithmicdo}
\algrenewtext{ForAll}[2] {\algorithmicforall\ #1 \algorithmicin\ #2 \algorithmicdo}

\newenvironment{algo}[4]{
	\noindent\rule{\textwidth}{0.4pt}
	\begin{algorithmic}[1]
		\addtocounter{ALG@line}{-1}
		\Procedure{#1}{#2}
		\Require #3
		\Ensure #4
		\Statex }{
		\EndProcedure
	\end{algorithmic}
	\rule{\textwidth}{0.4pt}}

\title{Appunti di Algoritmi per l'ingegneria}
\author{Giacomo Simonetto}
\date{Secondo semestre 2025-26}

\begin{document}

\maketitle
\begin{abstract}
	Appunti del corso di Algoritmi per l'Ingegneria della facoltà di Ingegneria Informatica dell'Università di Padova.
\end{abstract}

\newpage

\tableofcontents

\newpage

\section{Cenni di storia dell'algoritmica}
\dots
\section{Formalizzazioni dei concetti base}
\subsection{Algoritmo}
\subsubsection*{Definizione formale}
Un algoritmo è una macchina a stati finiti che riceve una sequenza di input, effettua una sequenza finita di transizioni tra
gli stati interni della macchina e produce una sequenza di output.

Si osserva che tale definizione coincide con la definizione di una macchina di Turing.

\subsubsection*{Stopping problem}
Si definisce lo ``stopping problem'' come la proprietà di un algoritmo di non poter continuare all'infinito. Ad esempio non
esiste un algoritmo che calcoli tutte le cifre di $\pi$ o che verifica se un'equazione diofantea (o diofantina) ha soluzioni
intere.

\subsubsection*{Progettazione di un algoritmo}
La progettazione di un algoritmo è il processo di ideazione e formalizzazione di un algoritmo per risolvere uno specifico
problema computazionale. Questo processo coinvolge:
\begin{itemize}
	\item la conoscenza del dominio del problema ed eventuali proprietà che lo caratterizzano, in quanto potrebbero
	semplificare il processo di risoluzione
	\item la conoscenza di metodi generali di progettazione di algoritmi, come i paradigmi di programmazione
\end{itemize}

\subsubsection*{Analisi di un algoritmo}
L'analisi di un algoritmo è il processo di valutazione dell'algoritmo secondo i seguenti aspetti:
\begin{itemize}
	\item correttezza: valuta se l'algoritmo porta alla soluzione corretta del problema per ogni istanza
	\item efficienza: valuta il consumo di risorse in termini di passi effettuati (complessità temporale), di memoria
	occupata (complessità spaziale) ed eventuali altri parametri (es. consumo energetico)
\end{itemize}

\subsection{Problema computazionale}
\subsubsection*{Definizione formale}
Un problema computazionale è composto dai seguenti tre elementi:
\begin{itemize}
	\item un insieme \(\mathcal{I}\) di istanze, che rappresentano i possibili input del problema, si nota che alcuni
	problemi potrebbero avere un numero infinito di istanze, per cui \(\mathcal{I}\) potrebbe essere un insieme infinito
	\item un insieme \(\mathcal{S}\) di soluzioni, che rappresentano gli output del problema per ogni istanza
	di \(\mathcal{I}\)
	\item una relazione \(\pi \subseteq \mathcal{I} \times \mathcal{S}\) che associa ogni istanza \(i \in \mathcal{I}\)
	ad una o più soluzioni \(s_j \in \mathcal{S}\) valide per quella istanza (è un sottoinsieme del prodotto cartesiano
	tra \(\mathcal{I}\) e \(\mathcal{S}\))
\end{itemize}

\subsubsection*{Rappresentazione delle istanze e delle soluzioni come stringhe}
I dati di un'istanza e di una soluzione sono rappresentati come stringhe, ovvero sequenze finite di simboli. Si assumono
le seguenti convenzioni:
\[\begin{array}{c c l }
	1. & \quad \Sigma = \{ 0, \; 1, \; \dots \} & \qquad \text{l'alfabeto dei simboli delle stringhe} \\[6pt]
	2. & \quad \Sigma^0 = \{\varepsilon\} & \qquad \text{l'insieme della sola stringa vuota} \\[6pt]
	3. & \quad \Sigma^i & \qquad \text{l'insieme di tutte le stringhe di lunghezza} \; i \; \text{(potenza cartesiana)} \\[6pt]
	4. & \quad \Sigma^+ = \bigcup_{i = 1}^{\infty} \Sigma^i & \qquad \text{l'insieme di tutte le stringhe esclusa la stringa vuota} \\[6pt]
	5. & \quad \Sigma^* = \bigcup_{i =0}^{\infty} \Sigma^i & \qquad \text{l'insieme di tutte le stringhe inclusa la stringa vuota} \\[6pt]
	6. & \quad \mathcal{I} \subseteq \Sigma^*, \quad \mathcal{S} \subseteq \Sigma^* & \qquad \text{gli insiemi di istanze e soluzioni sono sottoinsiemi di } \Sigma^*
\end{array}\]
La taglia di una istanza è la quantità di informazione necessaria per rappresentarla ed è pari al numero di simboli che
compongono la stringa che la rappresenta, o meglio la sua lunghezza. Si noti che per un numero intero \(n\) rappresentato
in binario, la taglia è il numero di cifre che lo compongono \(\approx \log_{2} n\) bit.

\subsubsection*{Osservazioni sulla rappresentazione dei dati nell'insieme delle istanze}
L'insieme delle istanze \(\mathcal{I}\) impone il tipo di rappresentazione con cui i dati in input devono essere codificati.
Fornire i dati in una diversa rappresentazione potrebbe portare a soluzioni errate, ad esempio se l'algoritmo svolge la somma
di numeri in base 10 e gli vengono forniti in base 2.

Inoltre la scelta della rappresentazione influisce sull'efficienza dell'algoritmo in quanto alcune operazioni risultano
più efficienti in una rappresentazione rispetto ad un'altra. Ad esempio la convoluzione nel dominio del tempo è \(O(n^2)\),
mentre nel dominio della frequenza è \(O(n \log n)\).

All'interno dell'algoritmo è possibile aggiungere una fase di conversione dei dati da una rappresentazione ad un'altra
(ad esempio per ridurre la complessità computazionale di alcuni problemi), ma è sempre necessario specificare la
rappresentazione dei dati in ingresso.

\subsubsection*{Funzioni e algoritmi}
Il concetto di relazione è più generale di quello di funzione, in quanto una relazione può associare più elementi del
codominio (es. soluzioni) ad un unico elemento del dominio (es. un'istanza). Una funzione, invece, può associare al massimo
un elemento del codominio ad ogni elemento del dominio.

Un problema computazionale definisce una relazione tra l'insieme delle istanze e quello delle soluzioni in quanto una istanza
può avere più soluzioni valide. Un algoritmo, invece, definisce una funzione tra i due insiemi. A partire da una istanza,
restituirà soltanto una delle possibili soluzioni per tale istanza. È possibile, inoltre, che due algoritmi diversi risolvono
lo stesso problema computazionale, ma producano soluzioni diverse, entrambe valide, a partire dalla stessa istanza.

\subsubsection*{Analisi della correttezza di un algoritmo}
Un algoritmo \(\mathcal{A}\) risolve un problema computazionale \(\pi\) se e solo se, detta \(A\) la funzione di trasferimento
tra ingresso e uscita definita da \(\mathcal{A}\), valgono le seguenti condizioni:
\[1. \quad \mathcal{I}_\text{istanze} \subseteq I_\text{input} \qquad\quad 2. \quad \mathcal{S}_\text{soluzioni} \subseteq O_\text{output} \qquad\quad 3. \quad \forall i \in \mathcal{I} \text{ vale } (i, A(i)) \in \pi\]
\[\text{con} \qquad \pi \subseteq \mathcal{I}_\text{istanze} \times \mathcal{S}_\text{soluzioni} \qquad \text{e} \qquad A: I_\text{input} \to O_\text{output}\]
L'analisi della correttezza consiste in particolare nel verificare la terza condizione.

\subsubsection*{Analisi della complessità di un algoritmo}
L'analisi della complessità computazionale di un algoritmo consiste nel valutare il tempo necessario all'algoritmo per
raggiungere la soluzione a partire da una data istanza generica.

La misurazione del tempo effettivo impiegato dipende molto dall'architettura della macchina su cui avviene la misurazione.
Il valore ottenuto sarebbe troppo variabile e non permetterebbe di stimare in maniera corretta e generale la complessità
di un algoritmo. Per questo motivo si simula l'esecuzione su random access machines con instruction set semplici e tempo
di esecuzione costante per ogni istruzione. Proprio il tempo di esecuzione di una istruzione è usato come unità di misura
del tempo di esecuzione. Questa semplificazione permette di stimare la complessità di un algoritmo in maniera più
generale, misurando banalmente il numero di istruzioni eseguite, e permette inoltre di confrontare tra loro algoritmi diversi
prima di averli implementati su una macchina reale.

Si effettua l'analisi al caso peggiore (worst-case) definendo come tempo impiegato da un algoritmo \(\mathcal{A}\) per
risolvere un'istanza \(i\) di taglia \(n\) come:
\[T_\mathcal{A}(n) \; := \; \sup \;\; T_\mathcal{A}(x) \qquad \text{con} \; \left\{ x \; : \; |x| = n \;\; \wedge \;\; x \in \Sigma^n \right\}\]
L'analisi al caso migliore è poco utile, mentre l'analisi al caso medio è notevolmente più complessa da effettuare, in quanto
richiede di conoscere a priori la distribuzione di probabilità delle varie istanze per una specifica taglia \(n\). In alcuni
casi quest'ultima viene preferita all'analisi al caso peggiore, se il caso peggiore è molto raro e mediamente (per la maggior
parte dei casi) la complessità risulta inferiore, ad esempio per il quicksort o per l'algoritmo del simplesso.

\newpage

\subsection{Dimostrazione per induzione}
\subsubsection*{Assiomi di Peano dei numeri naturali}
\begin{enumerate}
	\item esiste un numero naturale \(0\)
	\item ogni numero naturale \(n\) ha un successore \(s(n)\)
	\item numeri diversi hanno successori diversi: \(n \neq m \implies s(n) \neq s(m)\)
	\item \(0\) non è il successore di nessun numero naturale: \(s(n) \neq 0 \;\; \forall n\)
	\item forma debole del principio di induzione:
	\[\left[ \raisebox{3pt}{ \(\underbrace{P(0)}_{\text{caso base}} \wedge \;\; \underbrace{\forall n \; (P(n) \implies P(n + 1))}_{\text{passo induttivo}} \)} \;\; \right] \implies P(n) \;\; \forall n\]
\end{enumerate}

\vspace{5pt}

\noindent
Il quinto assioma può essere sostituito con il principio del minimo intero che sostiene che ogni insieme ordinabile non
vuoto ha un elemento minimo. Il principio di induzione e quello del minimo intero sono equivalenti, siccome il primo
implica il secondo e viceversa.

\subsubsection*{Forma forte del principio di induzione}
Esiste una formulazione più forte del principio di induzione, equivalente alla forma debole enunciata sopra, ma che facilita
le dimostrazioni basate su questo principio.
\[ \left[ \raisebox{3pt}{ \(\underbrace{P(0)}_{\text{caso base}} \wedge \;\; \underbrace{\forall n \; ((P(0) \wedge P(1) \wedge \ldots \wedge P(n)) \implies P(n + 1))}_{\text{passo induttivo}} \)} \;\; \right] \implies P(n) \;\; \forall n\]
Tale forma forte è ad esempio usata per dimostrare la correttezza del quicksort e negli ordinali transfiniti.

\subsubsection*{Rafforzamento dell'ipotesi induttiva}
Alcune ipotesi induttive più ``permissive'' potrebbero non essere dimostrabili tramite il principio di induzione, anche se
sono vere. Si riescono invece a dimostrare ipotesi più forti e ambiziose che implicano quelle più deboli. Ad esempio
potrebbe non essere dimostrabile che \(T(n) \leq n\), ma si riesce a dimostrare che \(T(n) \leq n - 1\) che implica
\(T(n) \leq n\).

Questo fenomeno paradossale viene chiamato rafforzamento dell'ipotesi induttiva. Si spiega con il fatto che l'ipotesi
induttiva compare sia nelle premesse che nelle conseguenze e, un suo rafforzamento, potrebbe sia rafforzare le conseguenze
che rafforzare le premesse a sua volta. Alcune volte il rafforzamento delle premesse è maggiore rispetto a quello delle
conseguenze e ciò permette di dimostrare queste ultime più facilmente.

\subsubsection*{Mappatura sui numeri naturali}
Il principio di induzione è definito solo sui i numeri naturali e può essere applicato solo su proprietà che hanno come
argomenti numeri naturali. Per dimostrare la correttezza di un algoritmo su un certo insieme di input (che non necessariamente
coincide con l'insieme dei numeri naturali) è necessario effettuare una opportuna mappatura tra gli elementi di tale insieme
e i numeri naturali. Spesso si clusterizza l'insieme degli input in base alla loro taglia e si mappa ogni cluster ad un numero
naturale che corrisponde alla taglia degli input in quel cluster.

Ad esempio, si vorrebbe dimostrare una certa proprietà che vale per gli alberi binari completi, ma gli alberi binari completi
non sono mappabili 1:1 con i numeri naturali. Si clusterizzano, quindi, in base alla loro taglia (ad esempio in base al numero
di nodi interni o alla sua altezza) e si mappa ogni cluster ad un numero naturale equivalente proprio al numero di nodi interni
o alla sua altezza. Si applica quindi il principio di induzione dimostrando che la proprietà vale per ogni cluster.

\subsubsection*{Limiti dell'induzione}
Il principio di induzione si utilizza principalmente per dimostrare ipotesi formulate a priori. Non aiuta a formulare nuove
ipotesi, che vanno pensate con l'aiuto di altre proprietà o intuizioni.

\subsubsection*{Induzione parametrica}
L'induzione parametrica o induzione empirica è una tecnica che permette di individuare i parametri di una funzione generale
o di un modello generale che si ritiene soddisfare al meglio i dati. Può essere usata, ad esempio, nell'addestramento di
reti neurali. In pratica consiste nell'applicare il principio di induzione imponendo che la formula valga per il caso base
e per il passo induttivo.

Ad esempio di vuole calcolare la formula analitica della seguente sommatoria:
\[s(n) = \sum_{k=1}^{n} k^2\]
Si notano le seguenti proprietà:
\begin{itemize}
	\item \(s(n) \leq n^3\) se tutti i valori della sommatoria fossero \(n^2\)
	\item \(s(n) \geq n^3 / 8\) siccome almeno \(n / 2\) valori della sommatoria sono maggiori o uguali a \((n / 2)^2\)
	\item \(s(n) \approx n^3 / 3 - 1/3\) se si tratta la sommatoria come un'integrale finito
\end{itemize}
per cui si ipotizza che la formula analitica sia un'equazione di terzo grado:
\[s(n) = a n^3 + b n^2 + c n + d\]

\begin{itemize}
	\item \textbf{caso base} \\
	si impone che valga il caso base \(s(0) = 0\) da cui si ottiene \(d = 0\)
	\item \textbf{passo induttivo} \\
	si impone che valga il passo induttivo \(s(n + 1) = s(n) + (n + 1)^2\):
\begin{align*}
	a (n + 1)^3 + b (n + 1)^2 + c (n + 1) + d & = a n^3 + b n^2 + c n + d + (n + 1)^2 \\
	a n^3 + (3 a + b) n^2 + (3 a + 2 b + c) n + (a + b + c + d) & = a n^3 + b n^2 + c n + d + n^2 + 2 n + 1 \\
	3a n^2 + (3 a + 2 b) n + (a + b + c) & = n^2 + 2 n + 1
\end{align*}
da cui si ottengono le tre equazioni che determinano gli altri tre parametri \(a\), \(b\) e \(c\):
\[3a = 1 \implies a = 1/3 \qquad\quad 3a + 2b = 2 \implies b = 1/2 \qquad\quad a + b + c = 1 \implies c = 1/6\]
\end{itemize}

\noindent
Si osserva che è possibile trovare gli stessi quattro parametri imponendo che la formula valga per quattro valori
particolari di \(n\) (ad esempio \(n = 0, 1, 2, 3\)) e risolvendo il sistema di quattro equazioni in quattro incognite.
Tale procedimento, però, senza aver osservato all'inizio che la formula è un'equazione di terzo grado, non permette di
affermare che i parametri trovati valgano per ogni \(n\). L'equazione, ad esempio, potrebbe essere di grado superiore
per cui i quattro punti scelti non sarebbero stati sufficienti.

\subsection{Formule matematiche utili}
\vspace{5pt}
\[\begin{array}{c l}
	\displaystyle a \, ^{\log_b c} = c \, ^{\log_b a} & \quad \text{esponenziale di un logaritmo} \\[10pt]
	\displaystyle \sum_{k=0}^{n-1} r^k = \frac{r^n - 1}{r - 1} & \quad \text{serie geometrica di ragione \(r\)} \\[15pt]
	\begin{array}{c}
		\overbrace{z = a + ib}^{\text{cartesiana}} \;\; \leftrightarrow \;\; \overbrace{z = \rho (\cos \varphi + i \sin \varphi)}^{\text{trigonometrica o polare}} \;\; \leftrightarrow \;\; \overbrace{\; z = \rho \cdot e^{i \varphi} \;}^{\substack{\text{esponenziale} \\ \text{o di Eulero}}} \\[5pt]
		\begin{cases}
			\rho = \sqrt{a^2 + b^2} \\
			\varphi = \arctan(b / a)
		\end{cases} \longleftrightarrow \quad \begin{cases}
			a = \rho \cos(\varphi) \\
			b = \rho \sin(\varphi)
		\end{cases}
	\end{array} & \quad \begin{array}{c}
		\text{rappresentazioni di un numero} \\ \text{complesso con passaggio tra} \\ \text{le varie rappresentazioni}
	\end{array} \\[40pt]
	\displaystyle\sum_{k=0}^{n-1} \binom{n}{k} a^k b^{n-k} = (a + b)^n \qquad \sum_{k=0}^{n} \binom{n}{k} = 2^n & \quad \begin{array}{c}
		\text{binomio di Newton con} \\ \text{caso particolare per } a = b = 1 \\ \text{(sottosequenze di una stringa)}
	\end{array} \\[30pt]
\end{array}\]

\section{Paradigma Divide and Conquer}
\subsection{Proprietà di sottostruttura}
\subsubsection*{Idea generale}
Il paradigma divide and conquer (divide et impera) è un approccio che consiste nel suddividere un'istanza di grandi dimensioni
in una o più sottoistanze di dimensioni inferiori, più semplici da risolvere, per poi combinare le soluzioni di queste istanze ed
ottenere la soluzione dell'istanza originale.

\subsubsection*{Proprietà di sottostruttura}
Dato un problema \(\pi \subseteq \mathcal{I} \times \mathcal{S}\), il paradigma si compone di tre componenti:
\begin{itemize}[topsep=5pt, itemsep=3pt] 
	\item \(A_D : \mathcal{I} \to \mathcal{I}*\), funzione di divide che suddivide un'istanza \(i \in \mathcal{I}\) di
	taglia \(|i| > n_0\) in una sequenza di una o più sottoistanze più piccole \(\langle i_1, i_2, \dots i_k \rangle\)
	di taglia \(|i_j| < |i|\) per ogni \(j\)
	\item \(A_C : \mathcal{S}* \to \mathcal{S}\), funzione di combine che combina la sequenza di soluzioni
	\(\langle s_1, s_2, \dots s_k \rangle\) associate alle sottoistanze \(\langle i_1, i_2, \dots i_k \rangle\)
	nella soluzione \(s\) associata all'istanza \(i\)
	\item \(D\&C : \mathcal{I} \to \mathcal{S}\), chiamata ricorsiva che trasforma individualmente ogni sottoistanza \(i_j\)
	nella corrispettiva soluzione \(s_j\), eventualmente applicando le due funzioni \(A_D\) e \(A_C\) in modo ricorsivo
\end{itemize}

\noindent
Le due funzioni \(A_D\) e \(A_C\) devono soddisfare la seguente proprietà di sottostruttura:
\[\forall i \in \mathcal{I}, |i| > n_0 \text{ vale } \left[ \; D(i) = \; \langle i_1, \dots, i_k \rangle \; \wedge \; \forall j \; (i_j, s_j) \in \pi \; \wedge \; C(\langle s_1, \dots, s_k \rangle) = s \implies (i, s) \in \pi \; \right]\]

\subsubsection*{Chiamata ricorsiva (conquer) e definizione del caso base}
La chiamata ricorsiva \(D\&C(i)\) su un'istanza \(i\) è definita come segue:
\begin{itemize}
	\item se \(|i| > n_0\) allora \(D\&C(i) = A_C(D\&C^*(A_D(i)))\), ovvero riapplica il modello di divide and conquer
	\item se \(|i| \leq n_0\) allora utilizza una soluzione diretta ed immediata per risolvere \(i\)
\end{itemize}
In generale il parametro \(n_0\) è scelto in base al problema. Può essere il minimo valore per cui è possibile applicare
il modello di divide and conquer, oppure un valore sufficientemente grande sotto di cui si conosce un algoritmo più
efficiente per risolvere il problema.

\begin{algo}{D\&C}{i}{\text{istanza \(i \in \mathcal{I}\)}}{\text{soluzione \(s \in \mathcal{S}\)}}
\If{\(|i| \leq n_0\)} \Comment{caso base \(\qquad\qquad\qquad\qquad\qquad\qquad\qquad\qquad\qquad\qquad\;\)}
	\State \(s \gets A_\text{diretto}(i)\) \Comment{soluzione diretta ed immediata \(\quad\qquad\qquad\qquad\quad\,\)}
	\State \textbf{return} \(s\)
\Else \Comment{passo ricorsivo \(\qquad\qquad\qquad\qquad\qquad\qquad\qquad\qquad\qquad\)}
	\State \(\langle i_1, i_2, \dots i_k\rangle \; \gets A_D(i)\) \Comment{divide \(\quad\qquad\qquad\qquad\qquad\qquad\qquad\qquad\qquad\qquad\)}
	\For{\(j \gets 1\)}{\(k\)} \Comment{conquer \(\quad\qquad\qquad\qquad\qquad\qquad\qquad\qquad\qquad\quad\:\)}
		\State \(s_j \gets D\&C(i_j)\)
	\EndFor
	\State \(s \gets A_C(\langle s_1, s_2, \dots s_k \rangle)\) \Comment{combine \(\quad\qquad\qquad\qquad\qquad\qquad\qquad\qquad\qquad\quad\)}
	\State \textbf{return} \(s\)
\EndIf
\end{algo}

\newpage

\subsection{Analisi della correttezza}
\subsubsection*{Proprietà di sottostruttura in funzione di un parametro intero}
Per dimostrare la correttezza di un algoritmo divide and conquer, si utilizza la proprietà di sottostruttura e il principio
di induzione in forma forte. Per prima cosa si esprime la proprietà di sottostruttura in funzione di un parametro intero,
in modo da renderla più adatta alla dimostrazione induttiva, ovvero:
\[\forall i \in \mathcal{I} \; \underbrace{\, \left[ \; (i, \; D\&C(i)) \in \pi \; \right] \,}_{P(i)} \;\; = \;\; \forall n \in \mathbb{N} \; \underbrace{\, \left[ \forall i \in \mathcal{I} \text{ con } |i| = n \text{ vale } (i, D\&C(i)) \in \pi \; \right] \,}_{P(n)}\]

\subsubsection*{Dimostrazione con il principio di induzione in forma forte}
Si dimostra, infine, che \(P(n)\) è vera per ogni \(n \in \mathbb{N}\) tramite il principio di induzione in forma forte:
\begin{itemize}
	\item si dimostra che \(P(n)\) è vera per ogni \(n \leq n_0\) (caso base)
	\item si dimostra che se \(P(n)\) è vera per tutti i valori \(n \leq m\) allora anche \(P(m + 1)\) è vera (passo induttivo)
\end{itemize}

\subsection{Albero della ricorsione}
\subsubsection*{Albero della ricorsione}
L'albero della ricorsione è una rappresentazione grafica che delle chiamate ricorsive di un algoritmo divide and
conquer con le seguenti caratteristiche:
\begin{itemize}
	\item ogni nodo rappresenta una chiamata ricorsiva \(D\&C(i)\) per un'istanza \(i \in \mathcal{I}\)
	\item ogni nodo ha come figli i nodi che rappresentano le chiamate ricorsive \(D\&C(i_j)\) per ogni sottoistanza
	\(i_j\) ottenuta dalla procedura di divide \(A_D(i)\)
	\item la taglia dell'istanza dei nodi figli è inferiore alla taglia dell'istanza del nodo padre, in quanto la taglia
	dell'istanza diminuisce progressivamente ad ogni chiamata ricorsiva
	\item ogni foglia rappresenta una chiamata ricorsiva \(D\&C(i)\) per un'istanza \(i \in \mathcal{I}\) di taglia
	\(|i| \leq n_0\), ovvero un caso base
\end{itemize}

\subsubsection*{Esplorazione dell'albero della ricorsione}
L'esplorazione dell'albero della ricorsione avviene secondo la visita in-order, ovvero prima si visita il figlio sinistro,
poi il nodo padre, poi il figlio destro. In questo modo, dato il nodo dell'albero che si sta attualmente visitando, si è
certi che tutti i nodi a sinistra di esso sono stati visitati, e i nodi a destra di esso sono ancora da visitare.

La costruzione della soluzione avviene secondo la visita post-order, ovvero prima si risolvono le istanze dei figli, poi
viene costruita la soluzione per il nodo padre, secondo l'algoritmo DFS (Depth First Search).

Lo spazio occupato è proporzionale allo spazio richiesto da ciascuna chiamata ricorsiva e dal numero di chiamate ricorsive
attive contemporaneamente, ovvero alla profondità dell'albero della ricorsione.

\subsubsection*{Cenni sulla ricorsione nei primi linguaggi di programmazione}
La memoria a disposizione sui primi computer era molto limitata e non era sufficiente per implementare uno stack che isolasse
le variabili locali delle varie chiamate ricorsive attive. I primi linguaggi di programmazione, infatti, salvavano le variabili
locali in posizioni fisse di memoria, comuni ad ogni chiamata ricorsiva. Per questo motivo, all'inizio la ricorsione non era
supportata. Il primo linguaggio di programmazione a supportare nativamente la ricorsione è stato il linguaggio LISP, sviluppato
da John McCarthy nel 1958, ma il suo utilizzo rimase fortemente scoraggiato per anni.

Ad oggi, la ricorsione è supportata nativamente da tutti i principali linguaggi di programmazione ed esistono ottimizzazioni
a livello hardware (istruzioni macchina per la gestione dello stack) per renderla efficiente e non sprecare memoria.

\subsection{Analisi della complessità e formulone}
\subsubsection*{Equazione alle ricorrenze}
L'equazione alle ricorrenze è un'equazione che esprime la complessità temporale di un algoritmo divide and conquer
in funzione della complessità temporale delle chiamate ricorsive e delle operazioni di divide e combine. Di seguito
è riportata la forma generale con le proprietà indicate sotto:
\[T(n) \; = \; \begin{cases}
	\; T_0(n) & \text{se } n \leq n_0 \quad - \quad \text{caso base} \\[3pt]
	\; s(n) \cdot T(f(n)) + \omega(n) & \text{se } n > n_0 \quad - \quad \text{passo ricorsivo}
\end{cases}\]
\begin{itemize}
	\item \(T_0(n)\) è il costo dell'algoritmo diretto per risolvere un'istanza di taglia \(n \leq n_0\)
	\item \(s(n)\) è il numero di chiamate ricorsive effettuate per risolvere un'istanza di taglia \(n\) e coincide con
	il numero di nodi figli nell'albero della ricorsione
	\item \(f(n)\) è la taglia della sottoistanza, generata dalla procedura di divide eseguita su un'istanza di taglia
	\(n\), che viene passata alla chiamata ricorsiva
	\item \(T(f(n))\) è il costo di ogni chiamata ricorsiva per risolvere un'istanza di taglia \(f(n)\)
	\item \(\omega(n) = T_D(n) + T_C(n)\) è il costo complessivo delle procedure di divide e di combine
\end{itemize}

\vspace{5pt}

\noindent
Per ottenere il valore analitico ed eliminare le ricorrenze è necessario risolvere l'equazione attraverso un'analisi
analitica, o ricorrendo al formulone o al master theorem.

\subsubsection*{Iterate di \(f\)}
Si definisce l'iterata di una funzione \(f\) come la potenza della composizione di \(f\) con se stessa.
\[\begin{cases}
	f^{(0)}(n) = n \\[3pt]
	f^{(l)}(n) = f(f^{(l-1)})(n) & \text{per } l > 0
\end{cases}\]
Si osserva che \(f^{(i)}(n)\) corrisponde alla taglia dell'istanza all'i-esima chiamata ricorsiva. In particolare,
inizialmente \(f^{(i)}(n) = n\) per \(i = 0\) e progressivamente diminuisce ad ogni iterata, fino a raggiungere il caso
base \(f^{(k+1)}(n) \leq n_0\) per un determinato valore di \(i = k\). Si definisce \(k\) come il massimo valore di \(i\)
per cui la taglia dell'istanza è ancora maggiore di \(n_0\), ovvero:
\[k = f^*(n, n_0) := \max \left\{ k : f^{(k)}(n) > n_0 \right\}\]

\subsubsection*{Formulone}
Il formulone è una formula generale che risolve l'equazione generale alle ricorrenze indicata sopra.
\[
T(n) = \underbrace{
	\overbrace{\sum_{l=0}^{f^* (n, \, n_0)}}^{\begin{array}{c} \text{\scriptsize per ogni} \\[-3pt] \text{\scriptsize livello interno} \end{array}}
	\overbrace{\left[ \prod_{j=0}^{l-1} s \left(f^{(j)}(n) \right) \right. \!}^{\begin{array}{c} \text{\scriptsize \# nodi al} \\[-3pt] \text{\scriptsize livello \(l\)-esimo} \end{array}}
	\; \cdot \!\!\! \overbrace{\left. \phantom{\prod_{j}^{l} \!\!\!\!\!\!\!\!} \omega \left( f^{(l)}(n) \right) \right]}^{\begin{array}{c} \text{\scriptsize costo interno} \\[-3pt] \text{\scriptsize di ogni nodo} \\[-3pt] \text{\scriptsize al livello \(l\)-esimo} \end{array}}
	}_\text{lavoro compiuto nei nodi interni}
	\;\; + \;\;\;
	\underbrace{ \overbrace{\prod_{j=0}^{f^*(n, n_0)} s \left(f^{(j)}(n) \right) \!}^{\begin{array}{c} \text{\scriptsize \# di foglie} \end{array}}
	\; \cdot \; \overbrace{ \phantom{\prod_{j}^{f^*} \!\!\!\!\!\!\!\!} T_0 \left( f^{(f^*(n, n_0) + 1)}(n) \right)}^{\begin{array}{c} \text{\scriptsize costo di ogni foglia}\end{array}}
	}_\text{lavoro compiuto nelle foglie}
\]

\vspace{7pt}

\noindent
Il formulone introduce le seguenti piccole semplificazioni che devono essere soddisfatte affinché la formula
indicata sopra sia valida:
\begin{itemize}
	\item il numero di figli dei nodi di un certo livello è costante, ovvero \(s(f^{(l)}(n)) = s_l\) per ogni livello \(l\)
	\item la taglia di tutti i nodi di un certo livello è costante, ovvero \(f^{(l)}(n) = f_l\) per ogni livello \(l\)
\end{itemize}

\subsubsection*{Formulone applicato ad un caso particolare}
Di seguito un esempio di applicazione del formulone ad un caso particolare con:
\begin{itemize}
	\item \(s(n) = a\): numero di figli costante per ogni nodo interno
	\item \(f(n) = n/2\): la taglia dell'istanza si dimezza ad ogni chiamata ricorsiva
	\item \(\omega(n) = \beta \, n\): costo di divide and conquer lineare per ogni nodo interno
	\item \(T_0(n) = 1\): costo costante per ogni foglia
\end{itemize}
\begin{align*}
	T(n) \;\; &= \sum_{l=0}^{(\log_2 n) - 1} \left( a^l \cdot \beta \; \frac{n}{2^l} \right) \; + \; a^{\log_2 n} \cdot T_0(n) \; = \; \beta \, n \cdot \! \sum_{l=0}^{(\log_2 n) - 1} \!\! \left( a/2 \right)^l \; + \; n^{\log_2 a} = \\[3pt]
	&= \;\; \beta \, n \cdot \frac{(a/2)^{\log_2 n} - 1}{(a/2) - 1} + n^{\log_2 a} \; = \; \frac{\beta}{a/2 - 1} \cdot n \left( n^{\log_2 a/2} - 1\right) + n^{\log_2 a} = \\[3pt]
	&= \; \underbrace{\frac{\beta}{a/2 - 1} \cdot (n^{\log_2 a} - n)}_{\text{costo dei nodi interni}} \;\; + \!\! \underbrace{ \;\; n^{\log_2 a} \;\;}_{\text{costo delle foglie}}
	\!\! = \;\; \Theta \left( n^{\log_2 a} \right)
\end{align*}

\subsubsection*{Complessità degli algoritmi di ordinamento basati sui confronti}
Al momento non si conoscono funzioni che legano problemi alla relativa complessità asintotica e più in generale non si
conoscono metodi per individuare in maniera assoluta i lower bound di problemi, ad eccezione di quelli deducibili banalmente.

Ad esempio gli algoritmi di ordinamento basati sui confronti hanno lower bound banale di \(\Theta(n \log n)\), siccome
\(n \log n\) è la quantità di informazione necessaria a rappresentare tutte le possibili permutazioni di \(n\) elementi ed
equivale al numero minimo di confronti (rappresentabili in bit) necessari per distinguere tutte le possibili permutazioni.
\[\# \; \text{permutazioni} = n! \approx n^n \qquad\quad \text{informazione}(\text{perm} \,_i) = \log_2(n!) \approx \log_2(n^n) = n \log_2 n\]

\subsection{Master theorem}
\subsubsection*{Assunzioni iniziali}
Il master theorem è un teorema che fornisce una soluzione asintotica immediata in funzione di due parametri \(a\) e \(b\)
per una classe di equazioni alle ricorrenze con le seguenti caratteristiche:
\begin{itemize}
	\item \(s(n) = a\): numero di figli costante per ogni nodo interno
	\item \(f(n) = n/b\): la taglia dell'istanza si divide per una costante \(b > 1\) ad ogni chiamata ricorsiva
	\item \(n_0 = 1\): caso base per istanze di taglia 1
\end{itemize}

\[T(n) \; = \; \begin{cases}
	\; T_0(n) & \text{se } n \leq 1 \\[3pt]
	\; a \cdot T(n/b) + \omega(n) & \text{se } n > 1
\end{cases}\]

\subsubsection*{Funzione soglia}
Si definisce una \textbf{funzione soglia} \(\sigma(n)\) che rappresenta il numero di foglie dell'albero della ricorsione.
Tale funzione permette di calcolare in maniera immediata il costo complessivo dovuto al lavoro compiuto nelle foglie che
funge da lower bound della complessità dell'algoritmo.
\[\sigma(n) = n^{\log_b a} \qquad T(n) \, \geq \, T_{\text{foglie}}(n) \, = \, T_0(n) \cdot \# \; \text{foglie} \, = \, T_0(n) \cdot \sigma(n)\]

\subsubsection*{Classi di funzioni e relative complessità asintotiche}
Paragonando il costo del lavoro compiuto da ogni nodo interno \(\omega(n)\) con la funzione soglia \(\sigma(n)\), si
individuano tre casi per cui è possibile calcolare la complessità asintotica di \(T(n)\) in maniera immediata:
\[T(n) \; = \; \begin{cases}
	\; \displaystyle \Theta (n^{\log_b a}) & \text{\footnotesize se il lavoro nei nodi interni \(\omega(n)\) è trascurabile rispetto a quello delle foglie \(\sigma(n)\)} \\[3pt]
	\; \displaystyle \Theta (n^{\log_b a} \, \log_b n) & \text{\footnotesize se il lavoro nei nodi interni \(\omega(n)\) è dello stesso ordine di quello delle foglie \(\sigma(n)\)} \\[3pt]
	\; \displaystyle \Theta (\omega(n)) & \text{\footnotesize se il lavoro nei nodi interni \(\omega(n)\) domina quello delle foglie \(\sigma(n)\)}
\end{cases}\]

\subsubsection*{Analisi del caso 1}
Si verifica quando il lavoro nei nodi interni \(\omega(n)\) è trascurabile rispetto a quello delle foglie \(\sigma(n)\):
\[\exists \varepsilon > 0 \text{ tale che } \omega(n) = O \! \left(n^{(\log_b a )- \varepsilon}\right) = \sigma(n) \cdot O \! \left(n^{-\varepsilon}\right)\]
Di seguito si osserva che calcolando il contributo complessivo dei nodi interni usando \(\omega(n) = n ^{(\log_b a) - \varepsilon}\),
il risultato è comparabile con quello delle foglie e non fa aumentare la complessità asintotica di \(\Theta(n^{\log_b a})\).
\begin{align*}
	T_{\text{int}}(n) &= \sum_{l=0}^{(\log_b n) - 1} a^l \cdot \omega \! \left( \frac{n}{b^l} \right) \; = \; \sum_{l=0}^{(\log_b n) - 1} a^l \cdot \left( \frac{n}{b^l} \right)^{(\log_b a) - \varepsilon} \; = \; \frac{n^{\log_b a}}{n^{\varepsilon}} \; \cdot \!\! \sum_{l=0}^{(\log_b n) - 1} \frac{a^l}{b^{l (\log_b a) - l \varepsilon}} \; = \\[3pt]
	&= \;\; \frac{n^{\log_b a}}{n^{\varepsilon}} \; \cdot \!\! \sum_{l=0}^{(\log_b n) - 1} \frac{a^l}{a^l \cdot b^{-l \varepsilon}} \; = \; \frac{n^{\log_b a}}{n^{\varepsilon}} \; \cdot \!\! \sum_{l=0}^{(\log_b n) - 1} {(b^{\varepsilon})}^l \; = \; \frac{n^{\log_b a}}{n^{\varepsilon}} \cdot \frac{(b^{\varepsilon})^{(\log_b n)} - 1}{b^{\varepsilon} - 1} \; = \\[3pt]
	&= \;\; \frac{n^{\log_b a}}{n^{\varepsilon}} \cdot \frac{n^{\varepsilon} - 1}{b^{\varepsilon} - 1} \; \approx \; \frac{n^{\log_b a}}{n^{\varepsilon}} \cdot \frac{n^{\varepsilon}}{b^{\varepsilon} - 1} \; = \; \frac{n^{\log_b a}}{b^{\varepsilon} - 1} \; = \; \frac{\sigma(n)}{b^{\varepsilon} - 1} \; = \; \Theta(\sigma(n))
\end{align*}
Si osserva che per \(\varepsilon \to 0^+\) il coefficiente moltiplicativo è molto elevato, ma comunque costante, per
cui asintoticamente trascurabile.

\subsubsection*{Analisi del caso 2}
Si verifica quando il lavoro nei nodi interni \(\omega(n)\) è dello stesso ordine di quello delle foglie \(\sigma(n)\):
\[\omega(n) = \Theta(n^{\log_b a}) = \Theta (\sigma(n))\]
Di seguito si dimostra che il contributo complessivo dei nodi interni ha una complessità asintotica maggiore di quella delle
foglie e domina la complessità totale dell'algoritmo, che risulta essere \(\Theta(n^{\log_b a} \, \log_b n)\).
\begin{align*}
	T_{\text{int}}(n) &= \sum_{l=0}^{(\log_b n) - 1} \left( a^l \cdot \omega \! \left( \frac{n}{b^l} \right) \right) \; = \; \sum_{l=0}^{(\log_b n) - 1} \left( a^l \cdot \left( \frac{n}{b^l} \right)^{\log_b a} \right) \; = \; n^{\log_b a} \; \cdot \!\! \sum_{l=0}^{(\log_b n) - 1} \left( \frac{a^l}{b^{l (\log_b a)}} \right) \\[3pt]
	&= n^{\log_b a} \; \cdot \!\! \sum_{l=0}^{(\log_b n) - 1} \frac{a^l}{a^l} \; = \; n^{\log_b a} \cdot (\log_b n) \; = \; \Theta(n^{\log_b a} \, \log_b n)
\end{align*}
Di seguito si osserva che il costo complessivo di ogni livello ha la stessa complessità asintotica di quello delle foglie,
per cui è possibile calcolare la complessità totale moltiplicando la complessità di ogni livello per il numero di livelli.
Gli algoritmi di questa tipologia con \(a = b\) (es. mergesort e quicksort) hanno la particolarità che il costo delle fasi di divide o
combine è proporzionale alla taglia dell'istanza.
\[T_{\text{\(l\)-esimo livello}}(n) = a^l \cdot \omega \! \left( \frac{n}{b^l} \right) = a^l \cdot \left( \frac{n}{b^l} \right)^{\log_b a} = a^l \cdot \frac{n^{\log_b a}}{b^{l (\log_b a)}} = a^l \cdot \frac{n^{\log_b a}}{a^l} = n^{\log_b a} = \sigma(n)\]

\newpage

\subsubsection*{Analisi del caso 3}
Si verifica quando il lavoro nei nodi interni \(\omega(n)\) domina quello delle foglie \(\sigma(n)\):
\[\exists \varepsilon > 0 \text{ tale che } \omega(n) = \Omega(n^{(\log_b a) + \varepsilon}) = \sigma(n) \cdot \Omega(n^{\varepsilon})\]
Dalla condizione precedente si può dedurre anche la seguente condizione di regolarità, ovvero che il lavoro interno complessivo
nei nodi figli (al livello inferiore) è minore del lavoro interno del nodo padre.
\[\exists c \in \mathbb{R} \text{ con } 0 < c < 1 \text{ tale che } \forall n \in \mathbb{N} \text{ si ha } a \cdot \omega \! \left( \frac{n}{b} \right) \leq c \cdot \omega(n)\]
Per la dimostrazione del terzo caso si utilizza la condizione di regolarità appena enunciata, in quanto rende lo svolgimento
più semplice e immediato. Di seguito sono riportati alcuni passaggi preliminari che verranno poi usati nella successiva
dimostrazione.
\[a\cdot \omega \! \left(\frac{n}{b}\right) \leq c\cdot \omega(n) \;\; \Rightarrow \;\; \omega \! \left(\frac{n}{b}\right) \leq \frac{c}{a} \cdot \omega(n) \;\; \Rightarrow \;\; \omega \! \left(\frac{n}{b^l}\right) \leq \left( \frac{c}{a} \right)^l \cdot \omega(n)\]
\[T_{\text{int}}(n) = \!\! \sum_{l=0}^{(\log_b n) - 1} a^l \cdot \omega \! \left( \frac{n}{b^l} \right) \; \leq \!\! \sum_{l=0}^{(\log_b n) - 1} a^l \cdot \frac{c^l}{a^l} \cdot \omega(n) \; = \; \omega(n) \cdot \!\!\! \sum_{l=0}^{(\log_b n) - 1} c^l \; < \; \omega(n) \cdot \frac{1}{1 - c} \; = \; \Theta(\omega(n))\]
Dall'analisi precedente si osserva che il costo complessivo dei nodi interni è proporzionale al costo di ogni nodo interno
e domina la complessità totale dell'algoritmo se il costo del singolo nodo interno è asintoticamente maggiore di quello delle
foglie, ovvero se \(\omega(n) = \Omega(n^{(\log_b a) + \varepsilon})\).

\subsubsection*{Considerazione sugli altri casi}
Esistono funzioni che non possono essere classificate in nessuno dei tre casi del master theorem, ad esempio quando
\(\omega(n) = n^{\log_b a} \log n\). In questi casi è necessario risolvere l'equazione alle ricorrenze attraverso
altri metodi risolutivi, ad esempio il formulone o l'analisi analitica.

\subsection{Applicazioni del divide and conquer}
Nei seguenti capitoli sono riportati alcuni esempi di applicazioni efficienti del paradigma divide and conquer:
\begin{itemize}
	\item esponenziazione con esponente intero
	\item moltiplicazione di polinomi tramite l'algoritmo di Karatsuba
	\item moltiplicazione di matrici tramite l'algoritmo di Strassen
	\item moltiplicazione di matrici supercircolanti tramite la trasformata di Hadamard
	\item moltiplicazione di polinomi tramite la trasformata di Fourier (FFT) e l'algoritmo di Cooley-Tukey
\end{itemize}

\newpage

\section{Esponenziazione con esponente intero}
\subsubsection*{Formalizzazione del problema}
Dato un elemento \(x\) proveniente da una struttura algebrica con operazione associativa (ad esempio numeri in \(\mathbb{N}\),
\(\mathbb{R}\), \(\mathbb{C}\), matrici, \dots) costante, si vuole calcolare \(x^n\) con \(n \in \mathbb{N}\). Siccome si
tratta \(x\) come costante, si esclude che possa essere un polinomio (in quanto aumenta di grado quando elevato a una potenza)
per semplificare la risoluzione.

\subsubsection*{Osservazioni sulla taglia del problema}
Essendo \(x\) costante, viene escluso dal calcolo della taglia che è determinata unicamente da \(n\). Si osserva che la
quantità di informazione necessaria a rappresentare \(n\) è \(\log_2 n\), per cui la taglia del problema in funzione di
cui va espressa la complessità è: \(\text{taglia}(n) = |n| = \log_2 n\), con \(n = 2 ^{\text{\, taglia}(n)}\).

\subsection{Algoritmo ovvio basato sulla definizione}
L'algoritmo ovvio per calcolare \(x^n\) si basa sulla definizione ricorsiva di potenza come composizione multipla
dell'operazione associativa della struttura algebrica, ovvero:
\[x^n := \begin{cases}
	\; x^{(0)} = 1 & \text{se } n = 0 \\[3pt]
	\; x^{(n-1)} \; \circ \; x & \text{se } n > 0 \quad - \quad \circ \;\; \text{è l'operazione associativa della struttura algebrica}
\end{cases}\]
La complessità è pari al costo di \(n - 1\) composizioni dell'operazione associativa, e risulta essere esponenziale in
funzione della taglia del problema, ovvero:
\[T(n) = n - 1 = 2^{\, \text{taglia}(n)} - 1 = \Theta \left( 2^{\, \text{taglia}(n)} \right)\]

\subsection{Algoritmo efficiente basato sul divide and conquer}
Si osserva che se l'operazione è associativa, come richiesto nelle ipotesi, è possibile riarrangiare le composizioni
di funzioni in modo da riutilizzare i risultati di composizioni già calcolate. Di seguito la definizione ricorsiva
dell'algoritmo efficiente basato sul divide and conquer:
\[x^n := \begin{cases}
	\; x^{(0)} = 1 & \text{se } n = 0 \\[3pt]
	\; x^{(1)} = x & \text{se } n = 1 \\[3pt]
	\; x^{(n/2)} \; \circ \; x^{(n/2)} & \text{se } n > 1 \text{ è pari} \\[3pt]
	\; x^{\lfloor n/2 \rfloor} \; \circ \; x^{\lfloor n/2 \rfloor} \; \circ \; x & \text{se } n > 1 \text{ è dispari}
\end{cases}\]
L'equazione alle ricorrenze per questo algoritmo al caso peggiore è la seguente, osservando che è sufficiente una chiamata
ricorsiva per calcolare \(x^{(n/2)}\) e, al caso peggiore, si aggiungono 2 operazioni di composizione.
\[T(n) = \begin{cases}
	\; 1 & \text{se } n = 0, 1 \\[3pt]
	\; T(n/2) + 2 & \text{se } n > 1
\end{cases}\]
Utilizzando il master theorem con \(a = 1\), \(b = 2\) e \(\omega(n) = O(1)\), si ottiene una complessità lineare
in funzione della taglia del problema, ovvero:
\[T(n) = \Theta \left( \log_2 n \right) = \Theta \left( \text{taglia}(n) \right)\]


\section{Moltiplicazione di polinomi - Karatsuba}
\subsubsection*{Definizione formale di polinomio e degree bound}
Un polinomio \(P(x)\) di grado \(N-1\) (con \(N\) termini) è definito come:
\[P(x) \; := \; \sum_{i=0}^{N-1} \, a_i \, x^i \qquad\qquad \text{con} \;\; a_0, a_1, \dotsc \; a_{N-1} \; \text{fissati}\]
Si osserva che un polinomio può essere visto anche come un numero espresso in base \(x\) con \(0\leq a_i < x\) secondo la
notazione posizionale.

Il degree bound \(\alpha\) è un limite superiore al grado di un polinomio e indica l'insieme di tutti i polinomi di
grado inferiore o uguale ad \(\alpha\), ovvero quelli con \(N - 1 \leq \alpha\).

\subsection{Prodotto tra polinomi usando la definizione}
\subsubsection*{Tecnica risolutiva}
Dati due polinomi \(P(x)\) e \(Q(x)\) di grado \(N - 1\), il loro prodotto \(P(x) \cdot Q(x)\) è anch'esso un poliomio
di grado \(2N - 2\), come illustrato di seguito.
\[P(x) \; := \; \sum_{i=0}^{N-1} \, a_i \, x^i \qquad\quad Q(x) \; := \; \sum_{i=0}^{N-1} \, b_i \, x^i \qquad\quad P(x) \cdot Q(x) \; = \; \sum_{i=0}^{N-1} \sum_{j=0}^{N-1} \, a_i \, b_j \, x^{i+j} \; = \; \sum_{k=0}^{2N-2} \, c_k \, x^k\]

\subsubsection*{Analisi della complessità}
Per calcolare il prodotto usando la definizione di prodotto tra polinomi sono necessarie:
\begin{itemize}
	\item \(N^2\) moltiplicazioni tra i coefficienti \(a_i\) e \(b_j\)
	\item \((N-1)^2 = N^2 - 1 - (2N - 2)\) addizioni tra i prodotti \(a_i \cdot b_j\) per ottenere i coefficienti \(c_k\)
\end{itemize}
Complessivamente il prodotto usando la definizione ha complessità \(\Theta(N^2)\).

\subsection{Prodotto tra polinomi con approccio divide and conquer}
\subsubsection*{Tecnica risolutiva}
I polinomi \(P(x)\) e \(Q(x)\), di grado \(N-1\), possono essere scomposti in polinomi più semplici di grado \(< N/2\)
e il prodotto tra \(P(x)\) e \(Q(x)\) è esprimibile in funzione del prodotto di questi polinomi più semplici, ovvero:
\[P(x) = A(x) + B(x) \cdot x^{N/2} \qquad\qquad Q(x) = C(x) + D(x) \cdot x^{N/2}\]
\[P(x) \cdot Q(x) = B(x) \, D(x) \cdot x^N + (A(x) \, D(x) + B(x) \, C(x)) \cdot x^{N/2} + A(x) \, C(x)\]

\subsubsection*{Equazione alle ricorrenze}
L'equazione alle ricorrenze per questo approccio è la seguente:
\[T(N) \; = \; \begin{cases}
	\; 1 & \text{se } N = 1 \\[3pt]
	\; 4 \cdot T(N/2) + \beta_{std} \; N & \text{se } N > 1
\end{cases}\]
\begin{itemize}
	\item il caso base è rappresentato dal prodotto di due polinomi di grado 0 che richiede un solo prodotto
	\item il passo ricorsivo prevede 4 chiamate ricorsive per calcolare i quattro prodotti \(AC\), \(AD\), \(BC\) e \(BD\)
	tra polinomi di grado \(< N/2\)
	\item il costo per ogni nodo interno è dato dalle somme dei coefficienti dei termini di grado uguale, che è proporzionale
	alla taglia \(N\)
\end{itemize}

\subsubsection*{Analisi della complessità}
Dal formulone si ottiene che la complessità rimane sempre \(\Theta(N^2)\), come evidenziato di seguito:
\begin{align*}
	T(N) \;\; &= \sum_{l=0}^{(\log_2 N) - 1} \!\! \left( 4^l \cdot \beta_{std} \; \frac{N}{2^l} \right) + 4^{\log_2 N} \cdot T_0(n) \; = \; \beta_{std} \; N \cdot \! \sum_{l=0}^{(\log_2 N) - 1} \!\! \left( 2^{2l}/2^l \right) + N^{\log_2 4} \; = \\[3pt]
	&= \;\; \beta_{std} \; N \cdot \sum_{l=0}^{(\log_2 N) - 1} \!\! 2^l + N^2 \; = \; \beta_{std} \; N \cdot \frac{2^{\log_2 N} - 1}{2 - 1} + N^2 \; = \\[3pt]
	&= \, \underbrace{\, \beta_{std} \; (N^2 - N) \,}_{\begin{array}{c} \text{\scriptsize costo delle somme} \\[-3pt] \text{ \scriptsize nei nodi interni} \end{array}}
	\;\; + \!\!\!\!\!\! \underbrace{\;\;\;N^2\;\;}_{\begin{array}{c} \text{\scriptsize costo dei prodotti} \\[-3pt] \text{\scriptsize nelle foglie} \end{array}} \!\!\!\!\!\! = \;\; \Theta \, (N^2)
\end{align*}

\subsection{Prodotto di polinomi con l'algoritmo di Karatsuba (1962)}
\subsubsection*{Tecnica risolutiva}
Karatsuba osserva che è possibile ricavare la somma dei due prodotti \(AD + BC\) riutilizzando i due prodotti \(AC\) e \(BD\)
già precedentemente calcolati come segue:
\[(A + B) \cdot (C + D) = AC + (AD + BC) + BD \implies AD + BC = (A + B) \cdot (C + D) - AC - BD\]
In questo modo il numero di prodotti da calcolare ricorsivamente si riduce da 4 a 3, con l'effetto di aumentare lievemente
il costo della fase di divide dovuto al calcolo delle somme \((A + B)\) e \((C + D)\).

\subsubsection*{Equazione alle ricorrenze ed analisi della complessità}
L'equazione alle ricorrenze per l'algoritmo di Karatsuba è la seguente:
\[T(N) \; = \; \begin{cases}
	\; 1 & \text{se } N = 1 \\[3pt]
	\; 3 \cdot T(N/2) + \beta_{kz} \; N & \text{se } N > 1
\end{cases}\]

\noindent
Dal formulone si ottiene che la complessità è \(\Theta(N^{\log_2 3}) \approx \Theta(N^{1.58})\), come evidenziato di seguito:
\begin{align*}
	T(N) \;\; &= \sum_{l=0}^{(\log_2 N) - 1} \!\! \left( 3^l \cdot \beta_{kz} \; N / 2^l \right) + 3^{\log_2 N} \cdot T(1) \; = \; \beta_{kz} \; N \cdot \! \sum_{l=0}^{(\log_2 N) - 1} \!\! \left( (3/2)^l \right) + N^{\log_2 3} \; = \\[3pt]
	&= \;\; \beta_{kz} \; N \cdot \frac{(3/2)^{\log_2 N} - 1}{(3/2) - 1} + N^{\log_2 3} \; = \; 2 \beta_{kz} \; N \cdot (N^{\log_2 3 - \log_2 2} - 1) + N^{\log_2 3} \; = \\[3pt]
	&= \;\; \underbrace{\, 2 \beta_{kz} \; (N^{\log_2 3} - N) \,}_{\begin{array}{c} \text{\scriptsize costo delle somme} \\[-3pt] \text{\scriptsize nei nodi interni} \end{array}}
	\;\; + \!\!\!\!\!\! \underbrace{\;\;N^{\log_2 3}\;\;}_{\begin{array}{c} \text{\scriptsize costo dei prodotti} \\[-3pt] \text{\scriptsize nelle foglie} \end{array}} \!\!\!\!\!\! = \;\; \Theta \, (N^{\log_2 3}) \; \approx \; \Theta \, (N^{1.58})
\end{align*}

\subsubsection*{Considerazioni sull'algoritmo di Karatsuba}
\begin{itemize}
	\item per valori di \(N\) piccoli, l'algoritmo standard (D\&C) rimane ancora più efficiente di quello di Karatsuba,
	in quanto il costo aggiuntivo della fase di divide è maggiore del risparmio ottenuto dalla chiamata ricorsiva in meno
	\item per valori di \(N\) molto grandi, è possibile migliorare ulteriormente l'efficienza asintotica nel calcolo del
	prodotto di polinomi utilizzando ad esempio l'algoritmo di Toom-Cook
	\item la vera novità dell'algoritmo di Karatsuba è stata quella di trovare il primo algoritmo sub-quadratico per
	risolvere un problema che tutti i più grandi matematici ritenevano impossibile da risolvere in meno di \(\Theta(N^2)\)
\end{itemize}

\section{Moltiplicazione di matrici - Strassen}
\subsubsection*{Note introduttive sulle matrici}
Per la seguente sezione relativa al prodotto di matrici, si considerano solo matrici quadrate \(n \times n\)
(\(\in \mathbb{R}^{n \times n}\)) con \(n = 2^k\) potenza di 2, per un certo \(k \in \mathbb{N}\). In questo modo
le matrici sono perfettamente scomponibili ricorsivamente in quattro sottomatrici di dimensione \(n/2 \times n/2\).

Per convenzione si assume che la taglia di una matrice è pari al numero di righe (o di colonne) di cui è composta,
ovvero \(A \in \mathbb{R}^{n \times n} \;\Rightarrow\; |A| = n\).

\subsection{Prodotto di matrici usando la definizione}
\subsubsection*{Tecnica risolutiva}
Date due matrici \(A, B \in \mathbb{R}^{n \times n}\), con \(n = 2^k\), il loro prodotto \(C = A \cdot B\) è sempre una
matrice \(C \in \mathbb{R}^{n \times n}\) i cui elementi \(c_{ij}\) sono calcolati come segue:
\[c_{ij} = \sum_{k=1}^{n} a_{ik} \cdot b_{kj} \qquad\qquad \text{per } i, j = 1, 2, \dotsc, n\]
In altre parole il prodotto di matrici è una matrice che ha come elementi il prodotto interno dei vettori riga e colonna
rispettivamente della prima matrice \(A\) e della seconda matrice \(B\).

\subsubsection*{Analisi della complessità}
Il prodotto interno tra un vettore riga e un vettore colonna richiede \(n\) moltiplicazioni e \(n-1\) addizioni, per un costo
totale di \(T(n) = 2n - 1\). Siccome vanno fatti \(n^2\) prodotti interni, il costo totale della moltiplicazione di matrici
risulta \(\Theta(n^3)\): \[T(n) = n^2 (2n - 1) = 2n^3 - n^2 = \Theta(n^3)\]

\subsection{Prodotto di matrici con approccio divide and conquer}
\subsubsection*{Tecnica risolutiva}
Le matrici \(A, B \in \mathbb{R}^{n \times n}\) possono essere scomposte ciascuna in quattro sottomatrici di dimensioni
\(n/2 \times n/2\), identificate secondo il seguente sistema di indici \(X_{11}, X_{12}, X_{21}, X_{22} \in \mathbb{R}^{n/2 \times n/2}\).
Il prodotto \(C = A \cdot B\) risulta essere in funzione dei prodotti tra queste sottomatrici.
\[A = \begin{bmatrix}
	A_{11} & A_{12} \\ A_{21} & A_{22}
\end{bmatrix} \qquad\qquad\quad B = \begin{bmatrix}
	B_{11} & B_{12} \\ B_{21} & B_{22}
\end{bmatrix} \qquad\qquad\quad C = \begin{bmatrix}
	C_{11} & C_{12} \\ C_{21} & C_{22}
\end{bmatrix}\]
\[\begin{array}{c c}
	C_{11} = A_{11} \cdot B_{11} + A_{12} \cdot B_{21} \qquad\quad & \qquad\quad C_{21} = A_{21} \cdot B_{11} + A_{22} \cdot B_{21} \\[5pt]
	C_{12} = A_{11} \cdot B_{12} + A_{12} \cdot B_{22} \qquad\quad & \qquad\quad C_{22} = A_{21} \cdot B_{12} + A_{22} \cdot B_{22}
\end{array}\]

\subsubsection*{Equazione alle ricorrenze}
L'equazione alle ricorrenze per questo approccio è la seguente:
\[T(n) \; = \; \begin{cases}
	\; 1 & \text{se } n = 1 \\[3pt]
	\; 8 \cdot T(n/2) + 4 n^2 & \text{se } n > 1
\end{cases}\]
\begin{itemize}
	\item il caso base è rappresentato dal prodotto di due matrici \(1 \times 1\) che richiede un solo prodotto
	\item il passo ricorsivo prevede 8 chiamate ricorsive per calcolare gli otto prodotti tra le sottomatrici di dimensione \(n/2 \times n/2\)
	\item il costo per ogni nodo interno è dato dalle 4 somme dei prodotti tra sottomatrici \(n/2 \times n/2\)
\end{itemize}

\subsubsection*{Analisi della complessità}
Dal formulone si ottiene che la complessità rimane sempre \(\Theta(n^3)\), come evidenziato di seguito:
\begin{align*}
	T(n) \;\; &= \sum_{l=0}^{(\log_2 n) - 1} \!\! \left( 8^l \cdot \left( \frac{n}{2^l} \right)^2 \right) + 8^{\log_2 n} \cdot T_0(n) \; = \; n^2 \cdot \! \sum_{l=0}^{(\log_2 n) - 1} \!\! 2^l + n^{\log_2 8} = n^2 \cdot \frac{2^{\log_2 n} - 1}{2 - 1} + n^3 \; = \\[3pt]
	&= \underbrace{\;n^2 \cdot (n - 1)\;}_{\begin{array}{c} \text{\scriptsize costo delle somme} \\[-3pt] \text{\scriptsize dei nodi interni} \end{array}}
	+ \!\!\!\!\!\! \underbrace{\;\;\; n^3 \;\;}_{\begin{array}{c} \text{\scriptsize costo dei prodotti} \\[-3pt] \text{\scriptsize alle foglie} \end{array}}
	\!\!\!\!\!\! = \;\; 2n^3 - n^2 \; = \; \Theta(n^3)
\end{align*}

\subsubsection*{Vantaggi dell'algoritmo D\&C ricorsivo rispetto a quello iterativo basato sulla definizione}
Ad oggi le macchine moderne performano molto meglio con l'algoritmo D\&C ricorsivo rispetto a quello iterativo basato sulla
definizione in quanto si lavora con sottomatrici sempre più piccole fino a quando riescono ad essere totalmente memorizzate
all'interno della cache. In questo modo si riducono drasticamente i tempi di accesso alla memoria principale dovuti ai
continui cache miss quando si scorrono per intero righe e colonne per calcolare i prodotti interni.

\subsection{Prodotto di matrici con l'algoritmo di Strassen (1969)}
\subsubsection*{Tecnica risolutiva}
Ispirandosi a Karatsuba, Strassen osserva che è possibile ricavare gli 8 prodotti tra le sottomatrici \(A_{ij}\) e \(B_{ij}\)
riutilizzando solo 7 prodotti, calcolati su opportune somme e differenze di sottomatrici.
\vspace{-12pt}
\begin{center}
	\begin{minipage}{0.25\textwidth}
		\begin{align*}
			A_1 &= (A_{12} - A_{22}) \\
			A_2 &= (A_{11} + A_{22}) \\
			A_3 &= (A_{11} - A_{21}) \\
			A_4 &= (A_{11} + A_{12}) \\
			A_5 &= A_{11} \\
			A_6 &= A_{22} \\
			A_7 &= (A_{21} + A_{22})
		\end{align*}
	\end{minipage}
	\begin{minipage}{0.25\textwidth}
		\begin{align*}
			B_1 &= (B_{21} + B_{22}) \\
			B_2 &= (B_{11} + B_{22}) \\
			B_3 &= (B_{11} + B_{12}) \\
			B_4 &= B_{22} \\
			B_5 &= (B_{12} - B_{22}) \\
			B_6 &= (B_{21} - B_{11}) \\
			B_7 &= B_{11}
		\end{align*}
	\end{minipage}
	\begin{minipage}{0.4\textwidth}
		\begin{align*}
			P_i = A_i \cdot B_i \qquad \text{per } i = 1, 2, \dotsc, 7
		\end{align*}
		\vspace{-20pt}
		\begin{align*}
			C_{11} &= P_1 + P_2 - P_4 + P_6 \\
			C_{12} &= P_4 + P_5 \\
			C_{21} &= P_6 + P_7 \\
			C_{22} &= P_2 - P_3 + P_5 - P_7
		\end{align*}
	\end{minipage}
\end{center}

\subsubsection*{Osservazioni sui prodotti di Strassen}
\begin{itemize}
	\item se si sviluppano i conti per le quattro espressioni di \(C_{ij}\) si può notare che si ottengono esattamente
	le stesse espressioni dell'approccio D\&C standard
	\item esiste una simmetria tra le espressioni dei blocchi \(A_i\) e \(B_i\): tra \(A_1\) e \(B_5\), tra \(A_2\)
	e \(B_2\), tra \(-A_3\) e \(B_6\), tra \(A_4\) e \(B_3\), tra \(A_5\) e \(B_7\), tra \(A_6\) e \(B_4\), tra \(A_7\) e \(B_1\)
\end{itemize}

\subsubsection*{Equazione alle ricorrenze}
L'equazione alle ricorrenze per l'algoritmo di Strassen è la seguente:
\[T(n) \; = \; \begin{cases}
	\; 1 & \text{se } n = 1 \\[3pt]
	\; 7 \cdot T(n/2) + 18 \; n^2/4 & \text{se } n > 1
\end{cases}\]

\begin{itemize}
	\item il caso base è rappresentato dal prodotto di due matrici \(1 \times 1\) che richiede un solo prodotto
	\item il passo ricorsivo prevede 7 chiamate ricorsive per calcolare i sette prodotti \(P_i\) tra i blocchi \(A_i\) e \(B_i\)
	di dimensione \(n/2 \times n/2\)
	\item il costo della fase di divide consiste nelle 10 somme e differenze tra sottomatrici \(A_{ij}, B_{ij} \in \mathbb{R}^{n/2 \times n/2}\)
	per calcolare i blocchi \(A_i\) e \(B_i\) (complessivamente \(10 \cdot n^2/4\))
	\item il costo della fase di conquer consiste nelle 8 somme e differenze tra i prodotti \(P_i \in \mathbb{R}^{n/2 \times n/2}\)
	per calcolare i blocchi \(C_{ij}\) (complessivamente \(8 \cdot n^2/4\))
\end{itemize}

\subsubsection*{Analisi della complessità}
Applicando il formulone si ottiene che la complessità è \(\Theta(n^{\log_2 7}) \approx \Theta(n^{2.81})\), come mostrato di seguito:
\begin{align*}
	T(n) \;\; &= \sum_{l=0}^{(\log_2 n) - 1} \!\! \left( 7^l \cdot \frac{18}{4} \cdot \left( \frac{n}{2^l} \right)^2 \right) + 7^{\log_2 n} \cdot T_0(n) \; = \; \frac{9}{2} \; n^2 \cdot \! \sum_{l=0}^{(\log_2 n) - 1} \!\! \left( (7/4)^l \right) + n^{\log_2 7} \; = \\[3pt]
	&= \;\; \frac{9}{2} \; n^2 \cdot \frac{(7/4)^{\log_2 n} - 1}{(7/4) - 1} + n^{\log_2 7} \; = \; \frac{9}{2} n^2 \cdot \frac{n^{\log_2 (7/4)} - 1}{3/4} + n^{\log_2 7} \; = \\[3pt]
	&= \;\; \frac{9}{2} \cdot \frac{4}{3} \; n^2 \cdot \left( n^{(\log_2 7) - 2} - 1 \right) + n^{\log_2 7} \; = \; 6 n^2 \cdot (n^{\log_2 7}/n^2 - 1) + n^{\log_2 7} \\[3pt]
	&= \; \underbrace{\; 6 n^{\log_2 7} - 6n^2 \;}_{\begin{array}{c} \text{\scriptsize costo delle somme} \\[-3pt] \text{\scriptsize dei nodi interni} \end{array}}
	\;\; + \!\!\!\! \underbrace{\;\; n^{\log_2 7} \;\;}_{\begin{array}{c} \text{\scriptsize costo dei prodotti} \\[-3pt] \text{\scriptsize alle foglie} \end{array}}
	\!\!\!\!\!\! = \;\; 7n^{\log_2 7} - 6n^2 \; = \; \Theta(n^{\log_2 7}) \; \approx \; \Theta(n^{2.81})
\end{align*}

\subsubsection*{Osservazioni generali sull'algoritmo di Strassen}
\begin{itemize}
	\item una volta le moltiplicazioni di numeri erano molto costose, ora i blocchi logici sono altamente ottimizzati; inoltre
	molti processori dispongono dell'istruzione MADD (multiply and add) che consente di accumulare il risultato di una
	moltiplicazione all'interno di un registro in un solo ciclo di clock
	\item se si volesse eseguire il quadrato di una matrice \(A\) con l'algoritmo di Strassen, si risparmierebbero 5 somme
	e sottrazioni siccome i blocchi \(A_i\) e \(B_j\) corrispondenti sarebbero uguali
	\item in generale non si può avere corrispondenza tra i blocchi \(A_i\) e \(B_i\), perché se così fosse, l'algoritmo
	diventerebbe commutativo portando sicuramente ad un risultato sbagliato siccome la moltiplicazione di matrici non è
	commutativa
	\item utilizzando la definizione, la complessità asintotica del prodotto di matrici era \(n^\Omega\) con \(\Omega = 3\),
	con Strassen si è riusciti a ridurre il valore di \(\Omega\) a \(\log_2 7 \approx 2.81\); ad oggi si conoscono algoritmi
	per il prodotto di matrici con \(\Omega \approx 2.37\), ma risultano poco vantaggiosi in quanto hanno costanti asintotiche
	molto elevate che li rendono svantaggiosi per valori di \(n\) piccoli
\end{itemize}

\subsection{Algoritmo ibrido per il prodotto di matrici}
\subsubsection*{Idea implementativa}
Per valori di \(n \leq n_0\) si osserva che l'algoritmo standard (D\&C) rimane ancora più efficiente di quello di Strassen.
Si vuole determinare tale valore \(n_0\) e sviluppare un algoritmo tale che per \(n > n_0\) utilizzi l'algoritmo di Strassen,
mentre per \(n < n_0\) utilizzi l'algoritmo standard (D\&C).

\subsubsection*{Calcolo del parametro \(n_0\)}
Per calcolare il parametro \(n_0\), si cerca il valore che minimizza la complessità dell'algoritmo di Strassen imponendo
che per \(n \leq n_0\) venga utilizzato l'algoritmo standard (D\&C), ovvero:
\begin{itemize}
	\item i livelli dell'albero della ricorsione arrivano a \((\log_2 (n/n_0)) - 1\)
	\item il costo delle foglie è pari alla complessità dell'algoritmo standard (D\&C), ovvero \(T_0(n) = 2n^3 - n^2\)
\end{itemize}
\begin{align*}
	T(n, n_0) \;\; &= \sum_{l=0}^{(\log_2 (n/n_0)) - 1} \!\! \left( 7^l \cdot \frac{18}{4} \cdot \left( \frac{n}{2^l} \right)^2 \right) \; + \; 7^{\log_2 (n/n_0)} \cdot \left( 2n_0^3 - n_0^2 \right) \; = \; \dots \; = \\[3pt]
	&= \;\; n^{\log_2 7} \cdot \underbrace{\frac{2 n_0 + 5}{{n_0}^{(\log_2 7) - 2}}}_{g(n_0)} - 6n^2
\end{align*}
Si procede a minimizzare la funzione \(T(n, n_0)\) per trovare il valore ottimale di \(n_0\):
\vspace{-5pt}
\[\min_{n_0} \, T(n, n_0) \;\; \longrightarrow \;\; \min_{n_0} \, g(n_0) \;\; \longrightarrow \;\; \begin{array}{l} n_{0, \min} \approx 10.477 \\[3pt] g(n_{0, \min}) = 3.895 \end{array}\]

\subsubsection*{Considerazioni sul parametro \(n_0\)}
Dal punto di vista teorico, si sceglie \(n_0 = 8\) in quanto è la potenza di 2 più vicina a \(n_{0, \min}\). Riscrivendo
l'equazione alle ricorrenze e la complessità asintotica dell'algoritmo ibrido, si ottiene:
\[T(n, n_{0, \min}) \; = \; \begin{cases}
	\; 2n^3 - n^2 & \text{se } n \leq 8 \\[3pt]
	\; 7 \cdot T(n/2) + 18 \; n^2/4 & \text{se } n > 8
\end{cases} \quad = \quad \begin{cases}
	\; 2n^3 - n^2 & \text{se } n \leq 8 \\[3pt]
	\; 3.918 \cdot n^{\log_2 7} - 6n^2 & \text{se } n > 8
\end{cases}\]

Dal punto di vista pratico, non è detto che \(n_0 = 8\) sia il valore ottimale, in quanto il modello teorico non tiene conto
di overhead dovuti al cambio di algoritmo, alla gestione della cache e altri fattori specifici della macchina su cui viene
eseguito l'algoritmo.

Ad oggi si sta sviluppando una branca di ricerca chiamata \textbf{adaptive software} che punta ad aggiustare dinamicamente i parametri
degli algoritmi di algebra lineare (come \(n_0\)) attraverso analisi e test pratici effettuati sulla macchina stessa, al fine
di ottenere le migliori prestazioni possibili per quella determinata macchina in uso.

\section{Moltiplicazione di matrici supercircolanti - Tr. di Hadamard}
\subsubsection*{Definizione di matrici supercircolanti}
Una matrice \(A \in \mathbb{C}^{n \times n}\), con \(n=2^k\) si dice supercircolante se:
\begin{itemize}
	\item per \(n = 1\) sempre (banalmente tutti i numeri sono considerati matrici supercircolanti)
	\item per \(n > 1\) la matrice si può dividere in quattro sottomatrici \(n/2 \times n/2\) uguali a due a due lungo la diagonale
	principale e la diagonale secondaria e tali sottomatrici sono a loro volta supercircolanti
	\[A = \begin{bmatrix} A_1 & A_2 \\ A_2 & A_1 \end{bmatrix} \]
\end{itemize}

\subsubsection*{Osservazioni sulle matrici supercircolanti}
\begin{itemize}
	\item nelle matrici supercircolanti, ogni riga è ottenuta tramite uno shift circolare della riga precedente
	\item le righe successive alla prima sono univocamente definite a partire dalla prima riga, detta riga generatrice,
	tramite opportuni shift circolari
	\item dato che è sufficiente la prima riga per definire l'intera matrice supercircolante, la taglia di una matrice
	supercircolante \(A\) è pari al numero di elementi della sua prima riga, ovvero \(n\)
	\item si osserva che l'insieme delle matrici supercircolanti è chiuso rispetto alla somma:
	\[A + B = \begin{bmatrix} A_1 & A_2 \\ A_2 & A_1 \end{bmatrix} + \begin{bmatrix} B_1 & B_2 \\ B_2 & B_1 \end{bmatrix} = \begin{bmatrix} A_1 + B_1 & A_2 + B_2 \\ A_2 + B_2 & A_1 + B_1 \end{bmatrix} \quad \text{supercircolante per def.}\]
	\item si osserva che l'insieme delle matrici supercircolanti è chiuso rispetto alla moltiplicazione:
	\[A \cdot B = \begin{bmatrix} A_1 & A_2 \\ A_2 & A_1 \end{bmatrix} \cdot \begin{bmatrix} B_1 & B_2 \\ B_2 & B_1 \end{bmatrix} = \begin{bmatrix} A_1 B_1 + A_2 B_2 & A_1 B_2 + A_2 B_1 \\ A_2 B_1 + A_1 B_2 & A_2 B_2 + A_1 B_1 \end{bmatrix} \quad \text{supercircolante per def.}\]
	\item siccome l'insieme delle matrici supercircolanti è chiuso rispetto alle operazioni di somma e moltiplicazione, allora
	è considerato un anello, in più si osserva che è anche commutativo, siccome le due operazioni godono della proprietà
	commutativa, ovvero \(A + B = B + A\) e \(A \cdot B = B \cdot A\)
\end{itemize}

\subsubsection*{Somme tra matrici supercircolanti}
Per sommare due matrici supercircolanti è sufficiente sommare le prime righe delle due matrici-addendi elemento per elemento, ottenendo
la prima riga della matrice risultante dalla somma. La complessità di questo approccio è \(T_{add}(n) = \Theta(n)\) con \(n\)
pari alla taglia delle matrici.

\subsection{Moltiplicazioni tra matrici supercircolanti}
\subsubsection*{Moltiplicazioni usando la definizione}
Usando la definizione di moltiplicazione di matrici, si osserva che si richiede il calcolo di 4 prodotti tra sottomatrici
supercircolanti di dimensione \(n/2 \times n/2\) e 2 somme tra matrici supercircolanti di dimensione \(n/2 \times n/2\).
L'equazione alle ricorrenze per questo approccio è la seguente:
\[T(n) \; = \; \begin{cases}
	\; 1 & \text{se } n = 1 \\[3pt]
	\; 4 \cdot T(n/2) + 2 \cdot n/2 & \text{se } n > 1
\end{cases} \;\; = \;\; \begin{cases}
	\; 1 & \text{se } n = 1 \\[3pt]
	\; 4 \cdot T(n/2) + n & \text{se } n > 1
\end{cases}\]
Dal master theorem con \(a = 4\), \(b = 2\) e \(\omega(n) = n\) si ottiene una funzione soglia di
\(\sigma(n) = n^{\log_2 4} = n^2\), tale per cui \(\omega(n) = \sigma(n) \cdot O(n^{-1})\), ovvero si ricade
nel primo caso del master theorem. La complessità è quindi \(T(n) = \Theta(\sigma(n)) = \Theta(n^2)\).

\newpage

\subsubsection*{Moltiplicazioni con approccio D\&C ottimizzato}
Si osserva che secondo opportune manipolazioni algebriche è possibile ricavare le due sottomatrici \(C_1\) e \(C_2\)
del prodotto attraverso soli 2 prodotti, con l'overhead di 6 somme e due moltiplicazioni per uno scalare.
\[\begin{array}{l}
	U = (A_1 + A_2) \cdot (B_1 + B_2) = A_1 B_1 + A_1 B_2 + A_2 B_1 + A_2 B_2 \\[3pt]
	V = (A_1 - A_2) \cdot (B_1 - B_2) = A_1 B_1 - A_1 B_2 - A_2 B_1 + A_2 B_2
\end{array} \qquad \begin{array}{l}
	C_1 = (U + V)/2 = A_1 B_1 + A_2 B_2 \\[3pt]
	C_2 = (U - V)/2 = A_1 B_2 + A_2 B_1
\end{array}\]
L'equazione alle ricorrenze per questo approccio è la seguente:
\[T(n) \; = \; \begin{cases}
	\; 1 & \text{se } n = 1 \\[3pt]
	\; 2 \cdot T(n/2) + 6 \cdot n/2 + 2 \cdot n/2 & \text{se } n > 1
\end{cases} \;\; = \;\; \begin{cases}
	\; 1 & \text{se } n = 1 \\[3pt]
	\; 2 \cdot T(n/2) + 4n & \text{se } n > 1
\end{cases}\]
Dal master theorem con \(a = 2\), \(b = 2\) e \(\omega(n) = 4n\) si ottiene una funzione soglia di
\(\sigma(n) = n^{\log_2 2} = n\) paragonabile a \(\omega(n)\). La complessità dal secondo caso del
master theorem risulta essere \(T(n) = \Theta(n \cdot \log n)\).

\subsection{Trasformata di Hadamard e Fast Hadamard Transform (FHT)}
\subsubsection*{Autovalori alle foglie dell'albero della ricorsione}
L'albero della ricorsione per il prodotto di matrici supercircolanti con approccio D\&C ottimizzato è un albero binario
completo in cui ogni nodo figlio sinistro riceve le somme tra le prime righe delle sottomatrici \(A_1\) e \(A_2\) e tra
quelle di \(B_1\) e \(B_2\), mentre ogni nodo figlio destro riceve le rispettive differenze, affinché siano moltiplicate opportunamente.

Le foglie ricevono due matrici degeneri in scalati \(a\) e \(b\) che è possibile dimostrare essere proprio gli autovalori
delle due matrici supercircolanti \(A\) e \(B\) di partenza.

\subsubsection*{Trasformata e prodotto di Hadamard}
È possibile interpretare le operazioni di divide come una trasformazione lineare che, partendo da una matrice supercircolante
rappresentata dalla sua prima riga, restituisce il vettore di autovalori della matrice supercircolante stessa. In questo modo
il prodotto di matrici supercircolanti può essere visto come un prodotto di autovalori, che risulta avere complessità
\(\Theta(n)\) con \(n\) pari alla taglia delle matrici supercircolanti (che è proprio il costo delle foglie). Una volta fatto
ciò, le operazioni di combine effettuano la trasformazione lineare inversa, restituendo la matrice supercircolante risultante
dal prodotto, rappresentata dalla sua prima riga.

La trasformazione lineare impiegata prende il nome di \textbf{trasformata di Hadamard}, mentre il prodotto membro a membro
tra i vettori di autovalori prende il nome di \textbf{prodotto di Hadamard}.

Come visto nell'analisi della complessità, la trasformata e la relativa antitrasformata di Hadamard vengono calcolate in
\(\Theta(n \cdot \log n)\), mentre il prodotto di Hadamard è calcolato in \(\Theta(n)\). Complessivamente, trasformata e
prodotto hanno complessità \(\Theta(n \cdot \log n)\) come l'approccio D\&C ottimizzato.

\subsubsection*{Cambio di rappresentazioni per risolvere problemi più efficientemente}
L'idea alla base del prodotto di matrici supercircolanti tramite la trasformata di Hadamard è quello di convertire i dati
in ingresso in una rappresentazione che permette di risolvere il problema richiesto in maniera più efficiente, per poi
convertire il risultato ottenuto nella rappresentazione originale. Altri esempi di algoritmi che adottano questa tecnica sono:
\begin{itemize}
	\item l'algoritmo di Strassen, anche se non è chiaro come interpretare la rappresentazione intermedia
	\item l'algoritmo di FFT, che utilizza la rappresentazione dei polinomi per punti ottenuta tramite la trasformata discreta
	di Fourier per calcolare il prodotto di polinomi in maniera più efficiente
\end{itemize}

\subsubsection*{Definizione di simmetria}
Una simmetria è una proprietà di invarianza rispetto ad una certa trasformazione. Ad esempio la proprietà commutativa è
l'invarianza rispetto alla trasformazione che scambia i due operandi di un'operazione.

Conoscere le simmetrie permette di identificare ed eliminare i gradi di libertà di un problema, riducendo sia il peso della
rappresentazione utilizzata che quello delle operazioni di calcolo.

\section{Moltiplicazione di polinomi e FFT - Fast Fourier Transform}
{
	\footnotesize Nota: la DFT (Discrete Fourier Transform) è la trasformata matematica discreta di Fourier, la FFT (Fast Fourier
	Transform) è l'algoritmo efficiente per calcolare la DFT.
}

\subsubsection*{Cenni sulla moltiplicazione di polinomi}
Dati due polinomi \(p(x)\) e \(q(x)\) rispettivamente di degree bound \(n\) e \(m\), il loro prodotto \(p(x) \cdot q(x)\)
è un polinomio di degree bound \(n + m\) il cui vettore di coefficienti \(c\) è dato dal prodotto di convoluzione tra i vettori
di coefficienti \(a\) e \(b\) dei polinomi \(p(x)\) e \(q(x)\), come illustrato di seguito:
\[p(x) = \sum_{i=0}^{n-1} a_i \, x^i = a_0 \, x^0 + a_1 \, x^1 + \dots + a_{n-1} \, x^{n-1} \qquad\quad q(x) = \sum_{i=0}^{m-1} b_i \, x^i = b_0 \, x^0 + b_1 \, x^1 + \dots + b_{m-1} \, x^{m-1}\]
\[p(x) \cdot q(x) = \sum_{i=0}^{n+m-2} c_i \, x^i \qquad\quad \text{con} \quad c_k = \!\! \sum_{\substack{i+j=k \\[2pt] 0 \leq i < n \\[2pt] 0 \leq j < m}} \!\! a_i b_j \qquad \text{per} \; k = 0, 1, \dotsc, n + m\]

\noindent
Applicando la pura definizione per il calcolo della convoluzione, la complessità risultante è quadratica nel degree bound dei
polinomi. La soluzione più efficiente che verrà descritta di seguito consiste in:
\begin{itemize}
	\item trasformare, tramite la DFT, i polinomi \(p(x)\) e \(q(x)\) dalla rappresentazione per vettori di coefficienti
	alla rappresentazione per vettori di punti o valori, valutando i polinomi nelle radici n-esime dell'unità, con complessità
	\(\Theta(n \cdot \log n)\) con \(n\) pari al degree bound dei polinomi
	\item eseguire il prodotto di Hadamard tra i vettori di punti o valori ottenuti, che equivale al prodotto di convoluzione
	dei vettori di coefficienti, con complessità lineare nel degree bound dei polinomi con complessità \(\Theta(n)\) con \(n\)
	pari al degree bound dei polinomi
	\item eseguire l'inversa della DFT per trasformare il vettore di punti o valori del polinomio risultante in un vettore
	di coefficienti con complessità \(\Theta(n \cdot \log n)\) con \(n\) pari al degree bound dei polinomi
\end{itemize}

\subsection{Rappresentazioni di un polinomio}
\subsubsection*{Rappresentazione tramite vettore di coefficienti}
Un polinomio \(p(x)\) di degree bound \(n\) può essere rappresentato tramite un vettore di coefficienti \(\underline{a}\)
di dimensione \(n\), in cui l'elemento \(a_i\) rappresenta il coefficiente del termine \(x^i\) del polinomio. Il prodotto
di due polinomi rappresentati in questo modo è dato dal prodotto di convoluzione tra i vettori di coefficienti.
\[\underline{a} = (a_0, a_1, \dotsc, a_{n-1}) \qquad \underline{b} = (b_0, b_1, \dotsc, b_{m-1}) \qquad \underline{c} = (c_0, c_1, \dotsc, c_{n+m-2})\qquad \underline{c} = \underline{a} * \underline{b}\]

\subsubsection*{Rappresentazione tramite vettore di punti o valori}
Un polinomio \(p(x)\) di degree bound \(n\) può essere rappresentato tramite un vettore di punti o valori \(\underline{y}\)
di dimensione \(n\), in cui l'elemento \(y_i\) rappresenta il valore del polinomio in un punto specifico \(x_i\). Esistono
infiniti modi di rappresentare un polinomio per punti, a seconda dei punti scelti, ma data una qualsiasi rappresentazione
esiste solo un polinomio di grado \(n\) associato a tale rappresentazione.

Il prodotto di due polinomi rappresentati per punti è dato dal prodotto element-wise tra i vettori di punti, avendo
l'accortezza di estenderli con nuovi elementi fino alla dimensione necessaria a rappresentare il polinomio risultante,
ovvero \(n + m - 1\) se i polinomi di partenza hanno degree bound \(n\) e \(m\).
\[(\underline{x}, \; \underline{y}) \text{ rappresentazione per punti di } p(x) \text{ con } \begin{array}{c}
	\underline{x} = (x_0, x_1, \dotsc, x_{n-1}) \\ \underline{y} = (y_0, y_1, \dotsc, y_{n-1})
\end{array} \text{ e } y_i = p(x_i)\]

\newpage

\subsection{Valutazione di un polinomio}
\subsubsection*{Valutazione di un polinomio usando la definizione}
Il processo di valutazione di un polinomio \(p(x)\) consiste nel calcolare il valore del polinomio in un punto specifico
\(x_i\). Tale processo permette di trasformare un polinomio dalla rappresentazione per vettore di coefficienti alla
rappresentazione per punti.

Il costo di valutazione di un polinomio \(p(x)\) di degree bound \(n\) in un punto \(x_i\) usando la definizione di polinomio
è lineare nel degree bound del polinomio.
\[T(n) \; = \underbrace{\; n - 2 \;}_{\text{potenze \(x^2, x^3, ...\)}} + \underbrace{\; n - 1 \;}_{\text{prodotti \(a_i \cdot x^i\)}} + \;\;\underbrace{\; n - 1 \;}_{\text{addizioni}} \; = \; 3n - 4 \; = \; \Theta(n)\]
Considerando che un polinomio deve essere valutato in \(n\) punti per essere trasformato in una rappresentazione per punti,
la complessità complessiva della trasformazione risulta quadratica.

\subsubsection*{Valutazione di un polinomio con l'algoritmo di Horner}
L'idea dell'algoritmo di Horner è quella di riscrivere un polinomio \(p(x)\) in una forma che permette di valutare
il polinomio ricorsivamente come segue:
\[p(x) = a_0 + a_1 x + a_2 x^2 + \dots + a_{n-1} x^{n-1} = a_0 + x \cdot \left( a_1 + a_2 x + \dots + a_{n-1} x^{n-2} \right)\]
L'equazione alle ricorrenze e la complessità dell'algoritmo di Horner per valutare il polinomio in un punto rimangono lineari
nel degree bound del polinomio, per cui non c'è un netto miglioramento rispetto all'approccio precedente.
\[T(n) \; = \; \begin{cases}
	\; 0 & \text{se } n = 1 \\[3pt]
	\; T(n-1) + 2 & \text{se } n > 1
\end{cases} \qquad\qquad T(n) = 2(n-1) = \Theta(n)\]

\subsubsection*{Valutazione di un polinomio sulle radici dell'unità}
Questo approccio si basa sulla possibilità di riutilizzare i conti fatti per valutare un polinomio \(p(x)\) in un punto
\(x_i\) per valutare lo stesso polinomio in altri punti \(x_j\), legati a \(x_i\) da una certa simmetria. Di seguito viene
illustrato un esempio per un polinomio di degree bound \(n = 8\).
\[\begin{array}{c c}
	p(x) \; = \; a_0 + a_1 x + a_2 x^2 + a_3 x^3 + a_4 x^4 + a_5 x^5 + a_6 x^6 + a_7 x^7 & \substack{\text{polinomio } p(x) \text{ di} \\ \text{degree bound } 8} \\[10pt]
	p(x) \; = \; \underbrace{(a_0 + a_2 x^2 + a_4 x^4 + a_6 x^6)}_{p^{[0]}(x^2)} \; + \; x \cdot \underbrace{(a_1 + a_3 x^2 + a_5 x^4 + a_7 x^6)}_{p^{[1]}(x^2)} & \substack{\text{si dividono le potenze} \\ \text{pari e dispari ottenendo} \\ \text{due polinomi in } x^2 \\ \text{di degree bound } 4} \\[20pt]
	p(x) = p^{[0]}(x^2) + x \cdot p^{[1]}(x^2) \qquad \begin{aligned} p(-x) &= p^{[0]}((-x)^2) - x \cdot p^{[1]}((-x)^2) \\ &= p^{[0]}(x^2) - x \cdot p^{[1]}(x^2) \end{aligned} & \substack{\text{le valutazioni del} \\\text{polinomio in } p(x) \text{ e} \\ p(-x) \text{ condividono gli} \\ \text{stessi due contribuiti} \\ p^{[0]}(x^2) \text{ e } p^{[1]}(x^2)} \\[22pt]
	\begin{array}{c}
		p^{[0]}(x^2) = p^{[0][0]}(x^4) + x^2 \cdot p^{[0][1]}(x^4) \\
		p^{[0]}(-x^2) = p^{[0][0]}(x^4) - x^2 \cdot p^{[0][1]}(x^4) \\[7pt]
		\begin{aligned}
			p(ix) &= p^{[0]}({ix}^2) + ix \cdot p^{[1]}({ix}^2) \\ &= p^{[0]}(-x^2) + ix \cdot p^{[1]}(-x^2)
		\end{aligned}
	\end{array} \begin{array}{c}
		p^{[1]}(x^2) = p^{[1][0]}(x^4) + x^2 \cdot p^{[1][1]}(x^4) \\
		p^{[1]}(-x^2) = p^{[1][0]}(x^4) - x^2 \cdot p^{[1][1]}(x^4) \\[7pt]
		\begin{aligned}
			p(ix) &= p^{[0]}({ix}^2) - ix \cdot p^{[1]}({ix}^2) \\ &= p^{[0]}(-x^2) - ix \cdot p^{[1]}(-x^2)
		\end{aligned}
	\end{array} & \substack{\text{calcolando i polinomi} \\ p^{[0]}(x^2) \text{ e } p^{[1]}(x^2) \\ \text{ si ottengono ``gratis''} \\ p^{[0]}(-x^2) \text{ e } p^{[1]}(-x^2) \\ \text{che forniscono anche} \\[1pt] \text{le due valutazioni} \\ p(ix) \text{ e } p(-ix) \\ \text{ con } i^2 = -1}
\end{array}\]
È possibile infine eseguire un'ulteriore passaggio ottenendo gli otto contributi \(p^{[i][j][k]}(x^8)\), ciascuno di
degree bound \(8/2/2/2 = 1\), che forniscono le valutazioni del polinomio \(p(x)\) anche per \(x \pm ix\) e \(-x \pm ix\).
Si valuta, \(p(x)\) sugli 8 punti \(\pm x, \; \pm ix,\; x \pm ix, \; -x \pm ix\) che sono proprio le radici ottave dell'unità.

Successivamente si dimostrerà che questo approccio (basato sulla DFT) ha complessità \(\Theta(n \cdot \log n)\) con \(n\)
pari al degree bound del polinomio, contro la complessità quadratica degli approcci precedenti.

\subsection{Interpolazione di un polinomio}
\subsubsection*{Interpolazione con la matrice di Vandermonde}
L'interpolazione di un polinomio \(p(x)\) consiste nel calcolare il vettore di coefficienti \(\underline{a}\) a partire
dal vettore di punti \(\underline{y}\). Tale processo permette di trasformare un polinomio dalla rappresentazione per
punti alla rappresentazione per vettore di coefficienti.

Un possibile approccio per l'interpolazione consiste nell'invertire la formula di valutazione di un polinomio di grado
\(n\), ottenendo così un sistema a \(n\) equazioni con \(n\) incognite (\(a_0, a_1, \dots a_{n-1}\)), i cui coefficienti
sono le potenze di \(x\).
\[p(x) = y = \sum_{i=0}^{n-1} a_i x^i \; \rightarrow \; \sum_{i=0}^{n-1} {x_j}^i a_i = y_j \; \rightarrow \; \underbrace{\begin{bmatrix}
	1 & x_0 & {x_0}^2 & \dots & {x_0}^{n-1} \\
	1 & x_1 & {x_1}^2 & \dots & {x_1}^{n-1} \\
	\vdots & \vdots & \vdots & \ddots & \vdots \\
	1 & x_{n-1} & {x_{n-1}}^2 & \dots & {x_{n-1}}^{n-1}
\end{bmatrix}}_{\text{matrice di Vandermonde} \;\; V(x_0, x_1, \dotsc, x_{n-1})} \cdot \begin{bmatrix}
	a_0 \\ a_1 \\ \vdots \\ a_{n-1}
\end{bmatrix} = \begin{bmatrix}
	y_0 \\ y_1 \\ \vdots \\ y_{n-1}
\end{bmatrix}\]
La matrice delle potenze di \(x\) è detta matrice di Vandermonde, denotata con \(V(x_0, x_1, \dots, x_{n-1})\). Il determinate
della matrice di Vandermonde è dato dalla seguente formula.
\[\det \left( V (\underline{x}) \right) \; = \prod_{0 \, \leq \, i \, < \, j \, \leq \, n-1} (x_j - x_i) \]
Se i punti \(x_0, x_1, \dotsc, x_{n-1}\) sono distinti, allora il determinante della matrice di Vandermonde è sempre
diverso da zero. Questo implica che la matrice è invertibile e il sistema di equazioni ha una soluzione unica. La complessità
rimane quadratica nel degree bound del polinomio dovuta al prodotto tra matrice inversa e vettore di punti, supponendo di
disporre dell'inversa della matrice di Vandermonde gratis.
\[\underline{a} = \left[ \; V(x_0, x_1, \dots x_{n-1}) \; \right]^{-1} \cdot \, \underline{y}\]

\subsubsection*{Interpolazione con il polinomio interpolatore di Lagrange}
Un altro approccio per l'interpolazione consiste nell'utilizzare il polinomio interpolatore di Lagrange
\[p(x) \; = \; \sum_{k=0}^{n-1} \; y_k \, L_k(x) \qquad\qquad L_i(x) \; = \; \frac{\displaystyle \prod_{j \neq i} (x - x_j)}{\displaystyle \prod_{j \neq i} (x_i - x_j)}\]
Si osserva che l'espressione in funzione di \(L_k(x)\) coincide con il polinomio \(p(x)\), in quanto, se tale polinomio
viene valutato in un punto \(x = x_k\) (con \(x_k \in\) rappresentazione per punti di \(p(x)\)), la sommatoria restituisce
proprio \(y_k\), che è proprio il valore del polinomio \(p(x)\) in \(x_k\), per ogni \(k \in [0, \, n-1]\), infatti:
\begin{itemize}
	\item \(L_i(x_k) = 1\) per \(i = k\), per la presenza di cancellazioni tra numeratore e denominatore
	\item \(L_i(x_k) = 0\) per ogni \(i \neq k\), siccome esisterà un fattore \(x_k - x_j = 0\) nel numeratore
	che non si semplificherà con nessun \(x_i - x_j\) del denominatore
\end{itemize}

\noindent
Il costo computazionale dell'approccio con il polinomio interpolatore di Lagrange rimane sempre quadratico nel degree bound
del polinomio da interpolare, in quanto è necessario calcolare \(n\) polinomi interpolatori di Lagrange \(L_k(x)\) (per
\(k \in [0, \, n-1]\)), ciascuno di grado \(n-1\).

\subsubsection*{Interpolazione di un polinomio sulle radici dell'unità tramite la matrice di Fourier}
Nei prossimi capitoli si osserverà che la valutazione di un polinomio \(p(x)\) sulle radici \(n\)-esime dell'unità
equivale al prodotto tra la matrice di Fourier e il vettore dei coefficienti del polinomio \(p(x)\). Di conseguenza,
per la trasformata inversa si utilizza l'inversa della matrice di Fourier e il vettore dei valori del polinomio \(p(x)\)
sulle radici \(n\)-esime dell'unità.

La valutazione (per come prima analizzato) costa \(\Theta(n \log_2 n)\), inoltre si dimostrerà che anche l'interpolazione costa \(\Theta(n \log_2 n)\),
ottenendo così un sistema efficiente per cambiare rappresentazione e utilizzare la rappresentazione per punti per il
calcolo del prodotto di polinomi.

\subsection{Matrice di Fourier e radici n-esime dell'unità}
\subsubsection*{Definizione di matrice di Fourier}
La matrice di Fourier \(F_n \in \mathbb{C}^{n \times n}\) è definita come la matrice di Vandermonde \(V\) i cui coefficienti
sono le radici \(n\)-esime dell'unità.
\[F_n = V\! \left( {\omega_n\!}^0, {\omega_n\!}^1, \dots, {\omega_n\!}^{n-1} \right) \qquad\qquad \text{con} \quad \omega_n = e^{-2\pi i / n} \;\; \text{radice principale dell'unità}\]

\subsubsection*{Radici n-esime dell'unità con relative proprietà e lemmi}
Una radice \(n\)-esima dell'unità è un numero complesso \(z_k = e^{i \varphi}\) con \(\varphi = 2\pi k/n\) e \(k \in [0, \; n\!-\!1]\):
\[{z_k}^n = [\exp (2 \pi i k / n)]^n = \exp (2 \pi i k / n \cdot n) = \exp (2 \pi i k) = 1\]
Una radice \(z_k\) è detta radice principale \(\omega_n\) se elevata per i vari valori di \(k \in [0, \; n\!-\!1]\)
genera tutte le restanti radici \(n\)-esime. Esistono più radici principali, solitamente si sceglie quella con \(k = 1\).
\[{\omega_n\!}^k = {z_1\!}^k = [\exp (2 \pi i / n)]^k = \exp (2 \pi i k / n) = z_k\]
Di seguito alcune proprietà e lemmi delle radici \(n\)-esime dell'unità:
\[\begin{array}{c c}
	{\omega_n}^i \cdot {\omega_n}^j = {\omega_n}^{i+j} & \qquad \begin{array}{c}\textbf{chiusura della moltiplicazione} \\ \text{l'insieme delle radici \(n\)-esime dell'unità} \\ \text{è chiuso rispetto alla moltiplicazione}\end{array} \\[15pt]
	{\omega_n}^k \cdot {\omega_n}^{n-k} = {\omega_n}^k \cdot {\omega_n}^{-k} = 1 & \qquad \begin{array}{c}\textbf{elemento inverso} \\ \text{l'insieme delle radici \(n\)-esime dell'unità possiede} \\ \text{l'elemento inverso } {\omega_n}^{n-k} \text{ rispetto alla moltiplicazione}\end{array} \\[15pt]
	{\omega_{d \cdot n}}^{d \cdot k} = {\omega_n}^k & \qquad \begin{array}{c}\textbf{lemma di cancellazione} \\ \text{moltiplicando per uno stesso fattore } d \text{ sia l'esponente} \\ \text{che il grado della radice, il risultato non cambia}\end{array} \\[15pt]
	{\omega_n}^{n/2} = {\omega_{2 \cdot n/2}}^{1 \cdot n/2} = \omega_2 & \qquad \begin{array}{c}\textbf{corollario del lemma di cancellazione} \end{array} \\[4pt]
	({\omega_n}^i)^2 = ({\omega_n}^{i + n/2})^2 = {\omega_{n/2}}^i & \qquad \begin{array}{c}\textbf{lemma di dimezzamento} \\ \text{le due radici } {\omega_n}^i \text{ e }{\omega_n}^{i + n/2} \text{ sono opposte} \\ \text{ed i loro quadrati sono uguali}\end{array} \\[15pt]
	\displaystyle \sum_{j=0}^{n-1} {({\omega_n}^i)}^j = \begin{cases} 0 & \text{se } i \; \text{mod} \; n = 0 \\[3pt] n & \text{altrimenti} \end{cases} & \qquad \begin{array}{c}\textbf{lemma di somma} \\ \text{la somma di tutte le potenze di una radice } n\text{-esima} \\ \text{dell'unità è zero se la radice è diversa da 1, altrimenti è \(n\)}\end{array}
\end{array}\]
\vspace{0pt}

\noindent
NOTA: il lemma di somma equivale alla somma degli elementi lungo la riga o la colonna \(i\)-esima della matrice di Fourier,
per cui la somma di tutte le righe e di tutte le colonne della matrice di Fourier è zero, tranne per la prima riga e la prima
colonna, la cui somma è \(n\).

\subsubsection*{Inversa della matrice di Fourier}
La matrice di Fourier è invertibile essendo una Vandermonde con determinante pari a \(n\) (per il lemma di somma) in quanto è
composta dalle radici \(n\)-esime dell'unità, tutte distinte tra loro.
\[F_{\omega_n}(i,j) = {\omega_n}^{ij} \qquad\qquad {F_{\omega_n}\!}^{-1}(i,j) = \frac{1}{n} \cdot {\omega_n}^{-ij} = \frac{1}{n} \cdot [F_{\omega_n}(i,j)]^{-1} = \frac{1}{n} \cdot F_{\omega_n}(n-i, \; j)\]
\vspace{-12pt}
\begin{align*}
	{F_n}^{-1} &= \frac{1}{n} \cdot V \!\left( \left({\omega_n}^{-1}\right)^0, \; \left({\omega_n}^{-1}\right)^1, \; \left({\omega_n}^{-1}\right)^2, \dots, \left({\omega_n}^{-1}\right)^{n-1} \right) \\
	&= \frac{1}{n} \cdot V \!\left( \left({\omega_n}^0\right)^{n-1}, \; \left({\omega_n}^1\right)^{n-1}, \; \left({\omega_n}^2\right)^{n-1}, \dots, \left({\omega_n}^{n-1}\right)^{n-1} \right) \\
	&= \frac{1}{n} \cdot V \!\left( {\omega_n}^0, \; {\omega_n}^{n-1}, \; {\omega_n}^{n-2}, \dots, {\omega_n}^1 \right)
\end{align*}

\noindent
La matrice inversa di Fourier è la matrice di Fourier scalata di \(1/n\) in cui le righe da \(1\) a \(n\!-\!1\) comprese
vengono invertite e l'insieme dei generatori è l'ultima riga della matrice di partenza.
\[({F_n\!}^{-\!1} \cdot F_n)(h, k) = \sum_{l=0}^{n-1} {F_n\!}^{-\!1}(h, l) \cdot F_n(l, k) = \sum_{l=0}^{n-1} \frac{1}{n} \, {\omega_n\!}^{-hl} \, {\omega_n\!}^{hk}= \frac{1}{n} \sum_{l=0}^{n-1} {\omega_n\!}^{l(k-h)} = \begin{cases} 1 & {\scriptstyle \text{se } h = k} \\ 0 & \text{\scriptsize altrimenti} \end{cases} = \; \mathbb{I}_n(h,k)\]

\subsection{Discrete Fourier Transform e algoritmo FFT}
\subsubsection*{Valutazione del polinomio sulle radici dell'unità con la matrice di Fourier}
Si osserva che il processo di valutazione di un polinomio \(p(x)\) sulle radici \(n\)-esime dell'unità, equivale al prodotto
tra la matrice di Fourier \(F_n\) e il vettore dei coefficienti \(\underline{a}\) del polinomio \(p(x)\). Come per la matrice
di Vandermonde, però, la complessità del prodotto applicando la pura definizione è \(\Theta(n^2)\).
\[\underline{y} \; = \begin{pmatrix}
	p(\omega_0) \\[3pt] p(\omega_1) \\[3pt] \vdots \\[3pt] p(\omega_{n-1})
\end{pmatrix} = \begin{pmatrix}
	\sum_{i=0}^{n-1} a_i \, {\omega_0}^i \\[3pt]
	\sum_{i=0}^{n-1} a_i \, {\omega_1}^i \\[3pt]
	\vdots \\[3pt]
	\sum_{i=0}^{n-1} a_i \, {\omega_{n-1}}^i
\end{pmatrix} =
\begin{pmatrix}
	{\omega_0}^0 & {\omega_0}^1 & \dots & {\omega_0}^{n-1} \\
	{\omega_1}^0 & {\omega_1}^1 & \dots & {\omega_1}^{n-1} \\
	\vdots & \vdots & \ddots & \vdots \\
	{\omega_{n-1}}^0 & {\omega_{n-1}}^1 & \dots & {\omega_{n-1}}^{n-1}
\end{pmatrix} \begin{pmatrix}
	a_0 \\ a_1 \\ \vdots \\ a_{n-1}
\end{pmatrix} = \; F_n \cdot \underline{a}\]

\subsubsection*{Formalizzazione della DFT}
Il prodotto tra la matrice di Fourier \(F_n\) e un generico vettore \(\underline{a}\) viene chiamato Discrete Fourier
Transform (DFT) del vettore \(\underline{a}\). È possibile, quindi, definire tale prodotto come una funzione come segue:
\[\underline{y} = F_n \cdot \underline{a} = \text{DFT}_n \left( \underline{a} \right)\]

\subsubsection*{Algoritmo efficiente per calcolare la DFT}
Si riprende la proprietà di sottostruttura individuata per valutare un polinomio \(p(x)\) per due generici punti \(x\) e
\(-x\) e la si applica per due generiche radici \(n\)-esime dell'unità opposte tra loro (\({\omega_n}^i, \; {\omega_n}^{i+n/2}\)).
\[\left\{ \; \begin{aligned}
	p(x) = p^{[0]}\!\left(x^2\right) + x \cdot p^{[1]}\!\left(x^2\right) \\
	p(-x) = p^{[0]}\!\left(x^2\right) - x \cdot p^{[1]}\!\left(x^2\right)
\end{aligned} \right. \quad \longrightarrow \quad \left\{ \; \begin{aligned}
	p \!\left({\omega_n}^i\right) = p^{[0]}\!\left({\omega_{n/2}}^i\right) + {\omega_n}^i \cdot p^{[1]}\!\left({\omega_{n/2}}^i\right) \\
	p \!\left({\omega_n}^{i + n/2}\right) = p^{[0]}\!\left({\omega_{n/2}}^i\right) - {\omega_n}^i \cdot p^{[1]}\!\left({\omega_{n/2}}^i\right)
\end{aligned} \right.\]
\[\text{con} \;\; \left({\omega_n}^i\right)^2 = \left({\omega_n}^{i + n/2}\right)^2 = {\omega_{n/2}}^i \;\; \text{per il lemma di dimezzamento o quadrato di radici opposte}\]

\noindent
Si esprime l'espressione sopra in funzione del vettore dei valori \(\underline{y}\) che assume il polinomio \(p(x)\) valutato
sulle radici \(n\)-esime dell'unità:
\[\underline{y} = \text{DFT}_n \left( \underline{a} \right) = \left\{ \;\; \begin{aligned}
	y_i = {y_i}^{[0]} + {\omega_n}^i \cdot {y_i}^{[1]} \\[4pt]
	y_{i + n/2} = {y_i}^{[0]} - {\omega_n}^i \cdot {y_i}^{[1]}
\end{aligned} \right. \quad \text{ per } i \in [0, \; n/2\!-\!1] \quad \text{con} \quad \begin{aligned}{}
	{y_i}^{[0]} &= p^{[0]} \! \left( {\omega_{n/2}}^i \right) \\ {y_i}^{[1]} &= p^{[1]} \! \left( {\omega_{n/2}}^i \right)
\end{aligned}\]

\noindent
Si osserva che i vettori \(\underline{y}^{[0]}\) e \(\underline{y}^{[1]}\) corrispondono alla valutazione dei polinomi
\(p^{[0]}(x)\) e \(p^{[1]}(x)\) sulle radici \(n/2\)-esime dell'unità, che possono essere calcolati riapplicando lo stesso
approccio in maniera ricorsiva.
\[y^{[0]} = F_{n/2} \cdot \underline{a}^{[0]} = \text{DFT}_{n/2} \! \left( \underline{a}^{[0]} \right) \qquad \text{e} \qquad y^{[1]} = F_{n/2} \cdot \underline{a}^{[1]} = \text{DFT}_{n/2} \! \left( \underline{a}^{[1]} \right)\]

\subsubsection*{Analisi della complessità dell'algoritmo efficiente per calcolare la DFT}
Il costo di ogni nodo interno è dato interamente dal costo del combine che comprende:
\begin{itemize}
	\item \(n/2\) moltiplicazioni tra i valori \(y_i^{[1]}\) e le radici \(n\)-esime dell'unità \({\omega_n}^i\)
	\item \(n/2\) addizioni per calcolare i valori \(y_i\)
	\item \(n/2\) sottrazioni per calcolare i valori \(y_{i + n/2}\)
	\item \(n/2 - 2\) prodotti per calcolare le radici \(n\)-esime dell'unità \({\omega_n}^i\) a partire da \(\omega_n\);
	il \(-2\) è dovuto al fatto che \({\omega_n}^0 = 1\) e \({\omega_n}^1 = \omega_n\) sono banali e non serve calcolarle
\end{itemize}
L'equazione alle ricorrenze e la complessità asintotica risultano, quindi, essere:
\[T(n) \; = \; \begin{cases}
	\; 0 & \text{se } n = 1 \\[3pt]
	\; 2 \cdot T(n/2) + 2n - 2 & \text{se } n > 1
\end{cases} \qquad\qquad T(n) = 2 n \log_2 n - 2 \log_2 n\]

\subsubsection*{Osservazione dei valori alle foglie dell'albero della ricorsione}
Costruendo l'albero della ricorsione si può osservare che i valori alle foglie sono ordinati secondo la permutazione
bit-reversal. Tale permutazione è ottenuta semplicemente invertendo l'ordine dei bit dell'indice \(i\) che identifica la foglia (es. \(i = 6\)
con rappresentazione binaria \(110\) diventa \(011\) che corrisponde a \(3\)).
\[\begin{array}{c c}
	\text{DFT}_{n}(\underline{a}) & (a_0, \; a_1, \; a_2, \; a_3, \; a_4, \; a_5, \; a_6, \; a_7) \\[3pt]
	\text{DFT}_{n/2}(\underline{a}^{[*]}) & (a_0, \; a_2, \; a_4, \; a_6) \quad\;\; (a_1, \; a_3, \; a_5, \; a_7) \\[3pt]
	\text{DFT}_{n/4}(\underline{a}^{[*][*]}) & (a_0, \, a_4) \quad (a_2, \, a_6) \quad\; (a_1, \, a_5) \quad (a_3, \, a_7) \\[3pt]
	\text{DFT}_{n/8}(\underline{a}^{[*][*][*]}) & (a_0) \;\; (a_4) \;\; (a_2) \;\; (a_6) \;\; (a_1) \;\; (a_5) \;\; (a_3) \quad (a_7)
\end{array}\]

\subsubsection*{Implementazione in pseudocodice dell'algoritmo FFT}
\begin{algo}{FFT}{$\underline{a}$}{rappresentazione tramite vettore di coefficienti $\underline{a}$ del polinomio $p(x)$}{rappresentazione tramite vettore di punti $\underline{y}$ del polinomio $p(x)$}
	\State \(n \gets \text{length}(\underline{a})\)
	\State
	\If{\(n = 1\)} \Comment{caso base}
		\State \textbf{return} \(a\)
	\EndIf

	\State
	\State \(\underline{a}^{[0]} \gets (a_0, a_2, a_4, \dots, a_{n-2})\) \Comment{fase di divide}
	\State \(\underline{a}^{[1]} \gets (a_1, a_3, a_5, \dots, a_{n-1})\)

	\State
	\State \(\underline{y}^{[0]} \gets \text{FFT}(\underline{a}^{[0]})\) \Comment{chiamate ricorsive}
	\State \(\underline{y}^{[1]} \gets \text{FFT}(\underline{a}^{[1]})\)

	\State
	\State \(\omega_n \gets \exp(2 \pi i /n)\) \Comment{fase di combine}
	\State \(\omega_i \gets 1\)
	\For{\(i \gets 0\)}{\(n/2-1\)}
		\State \(t \gets \omega_i \cdot y_i^{[1]}\)
		\State \(y_i \gets y_i^{[0]} + t\)
		\State \(y_{i + n/2} \gets y_i^{[0]} - t\)
		\State \(\omega_i \gets \omega_i \cdot \omega_n\)
	\EndFor
		
	\State \textbf{return} \(\underline{y}\)
\end{algo}

\newpage

\subsection{Moltiplicazione di polinomi con la DFT e padding dei coefficienti}
\subsubsection*{Padding dei coefficienti per il prodotto di polinomi}
Il prodotto di due polinomi \(p(x)\) e \(q(x)\) con degree bound rispettivamente di \(n\) e \(m\) è un polinomio
\(r(x) = p(x) \cdot q(x)\) con degree bound \(n + m - 1\). Per poter rappresentare correttamente tramite vettore di punti
il polinomio risultante è necessario valutare \(r(x)\) su almeno \(n + m - 1\) punti distinti.

Dal prodotto element-wise tra i vettori di punti di \(p(x)\) e \(q(x)\) valutati in \(\max\{n, m\}\) punti, si otterrà un
vettore di punti non sufficientemente lungo per rappresentare \(r(x)\). È necessario estendere le rappresentazioni tramite
vettori di coefficienti al degree bound \(n + m - 1\) del polinomio risultante, aggiungendo elementi nulli ai vettori di
coefficienti di \(p(x)\) e \(q(x)\). Questa operazione è detta padding dei coefficienti e gli elementi nulli aggiunti
corrispondo ai coefficienti delle potenze di grado superiore al degree bound dei due polinomi da moltiplicare.
\vspace{-5pt}
\begin{align*}
	\underline{a} = (a_0, a_1, \dots, a_{n-1}) \; \rightarrow \; \underline{a} = (a_0, a_1, \dots, a_{n-1}, \overbrace{0, 0, \dots, 0}^{m \; \text{elementi}}) \\
	\underline{b} = (b_0, b_1, \dots, b_{m-1}) \; \rightarrow \; \underline{b} = (b_0, b_1, \dots, b_{m-1}, \underbrace{0, 0, \dots, 0}_{n \; \text{elementi}}) \\
\end{align*}
\vspace{-33pt}

\subsubsection*{Padding per raggiungere la potenza di due più vicina}
Siccome l'algoritmo FFT è stato strutturato per funzionare con polinomi di degree bound pari ad una potenza di due,
è necessario estendere i vettori di coefficienti al degree bound pari alla potenza di due maggiore più vicina al degree
bound del polinomio risultante. Sarà necessario, quindi, effettuare un ulteriore padding e aggiungere ulteriori elementi
nulli ai vettori di coefficienti.
\begin{align*}
	\underline{a} = (a_0, a_1, \dots, a_{n-1}, 0, 0, \dots, 0) \; \rightarrow \; \underline{a} = \overbrace{(a_0, a_1, \dots, a_{n-1}, 0, 0, \dots, 0)}^{2^k \, \geq \, n + m - 1 \; \text{elementi}}) \\
	\underline{b} = (b_0, b_1, \dots, b_{m-1}, 0, 0, \dots, 0) \; \rightarrow \; \underline{b} = \underbrace{(b_0, b_1, \dots, b_{m-1}, 0, 0, \dots, 0)}_{2^k \, \geq \, n + m - 1 \; \text{elementi}}) \\
\end{align*}
\vspace{-33pt}

\subsubsection*{Note sulla complessità dopo il padding}
Definita \(N = 2^k\) la dimensione dei vettori di coefficienti dopo i due padding, la complessità del prodotto di polinomi
tramite la DFT è \(\Theta(N \log_2 N)\).

Supponendo che i degree bound dei polinomi \(p(x)\) e \(q(x)\) siano entrambi \(n\), il primo padding porterà la
dimensione dei vettori a \(2n\), pari al degree bound del polinomio risultante. La potenza più vicina a \(2n\) potrà
essere al massimo \(2 \cdot 2n = 4n\), siccome due potenze di due consecutive differiscono di un fattore 2.

Di conseguenza, la dimensione finale sarà al massimo \(N = 4n\), che provocherà un peggioramento di circa quattro volte
della complessità del prodotto, ma senza alterare quella asintotica di \(\Theta(n \log_2 n)\).

\subsubsection*{Note sull'interpretazione del padding in generale}
In generale non sempre è possibile applicare un padding all'input per ottenere una taglia di dimensione ottimale. Di seguito
alcune osservazioni utili per interpretare il padding in generale:
\begin{itemize}
	\item nel prodotto tra polinomi, il padding di elementi nulli non altera il polinomio, in quanto equivale ad aggiungere
	coefficienti nulli delle potenze di grado superiore
	\item nel mergesort, il padding potrebbe essere fatto con valori facili da discriminare (es. \(-\infty\) o \(+\infty\))
	per cui è possibile ottenere la stringa ordinata correttamente dopo averli eliminati dall'output
	\item nella trasformata di Fourier (e nella DCT), in generale, gli elementi da aggiungere al vettore di coefficienti
	affinché l'output non venga alterato sono molto complessi da determinare, per cui non è possibile applicare un padding
	in maniera efficiente
\end{itemize}

\noindent
NOTA: il padding di zeri nei vettori di input della DFT per la moltiplicazione di polinomi altera il risultato della DFT
perché banalmente si sta valutando il polinomio su più radici dell'unità. Tuttavia, in questo contesto non è un problema,
in quanto entrambe le rappresentazioni ottenute prima e dopo il padding (seppur sostanzialmente diverse) sono equivalenti
per rappresentare lo stesso polinomio di partenza.

\subsubsection*{Padding per \(k\)-potenza di un polinomio}
Si osserva che man mano che il polinomio viene elevato a potenza, il degree bound continua a crescere e sarebbe necessario
effettuare continui padding per garantire la valutazione su un numero di punti sufficientemente grande per rappresentare
i risultati intermedi e tale numero deve anche essere pari ad una potenza di 2.

In alternativa a tale processo di padding continuo, si può effettuare un padding iniziale già sufficientemente grande
per garantire che tutti i risultati intermedi e finali siano rappresentati correttamente.

Analizzando la complessità di entrambi i processi, si osserva che il costo complessivo dell'elevamento a potenza è
praticamente lo stesso ed è pari a \(\Theta(2 \, kn \log_2 kn)\), che non altera la complessità asintotica.

\subsection{Osservazioni utili per l'esame}
\subsubsection*{Vettore di output di una trasformazione data la matrice e un certo vettore di input}
Quando viene richiesto di calcolare il vettore di output di una trasformazione data la matrice relativa e un certo vettore
di input, è utile osservare che il vettore di output è dato dalla combinazione lineare delle colonne della matrice pesate
dai coefficienti del vettore di input.
\[\underline{y} = F \cdot \underline{a} = \sum_{i=0}^{n-1} F(*,i) \cdot a_i \qquad \text{con } F(*,i) \text{ colonna } i\text{-esima di } F\]
Nel caso della trasformata di Fourier esistono alcuni vettori notevoli:
\begin{itemize}
	\item \(\underline{a} = [1, 0, 0, \dots 0] \;\; \rightarrow \;\; \underline{y} = F_n \cdot \underline{a} = [1, 1, 1, \dots, 1]\)
	siccome la prima colonna della matrice di Fourier è composta da tutti 1
	\item \(\underline{a} = [0, 1, 0, 0, \dots 0] \;\; \rightarrow \;\; \underline{y} = F_n \cdot \underline{a} = [{\omega_n}^0, {\omega_n}^1, {\omega_n}^2, \dots, {\omega_n}^{n-1}]\)
	siccome la seconda colonna della matrice di Fourier è composta dalle radici \(n\)-esime dell'unità
	\item \(\underline{a} = [1, 1, 1, \dots 1] \;\; \rightarrow \;\; \underline{y} = F_n \cdot \underline{a} = [n, 0, 0, \dots, 0]\)
	siccome la somma degli elementi lungo le righe della matrice di Fourier è zero, tranne per la prima riga, la cui somma è \(n\) (lemma di somma)
	\item \(\underline{a} = (n,k)\text{-sparso}\) con \(n\) e \(k\) potenze di 2, è un vettore di \(n\) elementi con \(x_i = 0 \text{ per } i \text{ mod } k \neq 0\).
	Un caso che compare spesso sono i vettori \((8,4)\text{-sparsi} : \underline{a} = [p, 0, 0, 0, q, 0, 0, 0]\) la cui trasformata è:
	\(\underline{y} = F_8 \cdot [p, 0, 0, 0, q, 0, 0, 0] = p \cdot [1,1,1,1,1,1,1,1] + q \cdot [1,-1,1,-1,1,-1,1,-1]\)
\end{itemize}

\subsubsection*{Composizione di trasformate di Fourier}
La composizione di due trasformate di Fourier \({F_n}^2\) provoca uno scalamento di un fattore \(n\) e un ribaltamento degli
elementi del vettore di input, escluso il primo elemento, che rimane invariato.
\[{F_n}^2 \cdot \underline{a} = n \cdot (a_0, a_{n-1}, a_{n-2}, \dots, a_1) = n \cdot \underline{a}^*\]

\noindent
Siccome il ribaltamento applicato due volte riporta gli elementi del vettore di input alla posizione originale, la composizione
di quattro trasformate di Fourier \({F_n}^4\) provoca soltanto uno scalamento di \(n^2\) senza modificare la posizione degli
elementi del vettore di input.
\[{F_n}^4 \cdot \underline{a} = n^2 \cdot (a_0, a_1, a_2, \dots, a_{n-1}) = n^2 \cdot \underline{a}\]

\subsubsection*{Inversa della trasformata di Fourier}
Dalla composizione di due trasformate di Fourier si ottiene che l'inversa della trasformata di Fourier di un vettore
\(\underline{y}\) è data dalla trasformata di Fourier del vettore \(\underline{y}\) scalato di \(1/n\), a cui vengono
ribaltati gli elementi, escluso il primo.
%\[\begin{array}{c c c}
%	1. & \quad \underline{y} = F_n \cdot \underline{a} \;\; \rightarrow \;\; \underline{y} = \text{DFT}_n(\underline{a}) &
%	\quad \substack{\text{si riscrive la trasformata} \\ \text{come funzione matematica}} \\[7pt]
%	2. & \quad n \cdot \underline{a}^* = {F_n}^2 \cdot \underline{a} \;\; \rightarrow \;\; n \cdot \underline{a}^* = \text{DFT}_n(\text{DFT}_n(\underline{a})) &
%	\quad \substack{\text{dalla formula di composizione} \\ \text{di due trasformate di Fourier}} \\[7pt]
%	3. & \quad n \cdot \underline{a}^* = \text{DFT}_n(\text{DFT}_n(\underline{a})) = \text{DFT}_n(\underline{y}) &
%	\quad \substack{\text{dalla formula 1.}\; : \; \text{DFT}_n(\underline{a}) \; = \; \underline{y}} \\[7pt]
%	4. & \quad \underline{a}^* = \frac{1}{n} \cdot \text{DFT}_n(\underline{y}) &
%	\quad \substack{\text{invertendo la formula sopra}}
%\end{array}\]
\[n \cdot \underline{a}^* = {F_n}^2 \cdot \underline{a} = \text{DFT}_n(\text{DFT}_n(\underline{a})) = \text{DFT}_n(\underline{y}) \implies \underline{a}^* = \frac{1}{n} \cdot \text{DFT}_n(\underline{y}) \]
\[\text{con} \quad \underline{y} = \text{DFT}_n(\underline{a}) \quad \text{e} \quad \underline{a}^* = \left( a_0, a_{n-1}, a_{n-2}, \dots a_1 \right)\]

\section{Programmazione dinamica}
\subsection{Memoizzazione}
\subsubsection*{Concetto di memoizzazione per istanze ripetute}
Il paradigma della programmazione dinamica è un caso particolare di divide and conquer. Viene utilizzato nei problemi in cui
alcune sottoistanze vengono risolte più volte in punti diversi dell'albero della ricorsione. Per evitare il lavoro ridondante
e migliorare l'efficienza dell'algoritmo, la programmazione dinamica sfrutta la memoizzazione che consiste nel tenere traccia
delle varie sottoistanze risolte con le rispettive soluzioni, in modo da non doverle ricalcolare quando si presentano nuovamente.

La presenza di sottoistanze uguali è abbastanza rara, per cui il paradigma di programmazione dinamica viene impiegato
raramente e solo in determinati problemi con strutture particolari.

\subsubsection*{Approccio bottom-up iterativo e ordine di visita delle sottoistanze}
In alternativa all'approccio top-down ricorsivo tipicamente usato nel D\&C, è possibile risolvere i problemi di programmazione
dinamica con un approccio bottom-up iterativo. Durante questo processo, si esplorano tutte le sottoistanze dal caso base fino
a raggiungere l'istanza principale richiesta dall'utente e si memorizzano le soluzioni temporanee in una struttura dati, come
ad esempio una tabella.

Per fare ciò, in fase di progettazione, è necessario identificare un ordine di visita delle sottoistanze in modo tale da
garantire che una sottoistanza venga risolta solo dopo aver risolto tutte le sottoistanze da cui essa dipende. In altre
parole è necessario individuare un ordinamento topologico delle sottoistanze.

\subsection{Calcolo dell'\texorpdfstring{\(n\)}{n}-esimo numero di Fibonacci}
\subsubsection*{Formalizzazione del problema}
L'\(n\)-esimo numero di Fibonacci è definito ricorsivamente come segue:
\[F(n) = \begin{cases}
	n & \text{se } n = 0,1 \\
	F(n-1) + F(n-2) & \text{se } n > 1
\end{cases} \qquad \text{con} \qquad \mathcal{I} = \mathbb{Z}_{+,0} \qquad \mathcal{S} = \mathbb{Z}_{+,0}\]
L'istanza del problema è data da un numero \(n\), per cui la taglia risulta essere \(\log_2(n)\).

\subsubsection*{Soluzione ricorsiva non efficiente}
Applicando puramente la definizione ricorsiva illustrata sopra si ottiene il seguente algoritmo.

\begin{algo}{REC\_FIB}{n}{intero positivo $n \in \mathbb{Z}_{+,0}$}{$n$-esimo numero di Fibonacci $F(n) \in \mathbb{Z}_{+}$}
	\If{\(n = 0\) \textbf{or} \(n = 1\)}
		\State \Return \(n\)
	\Else
		\State \Return \(\text{REC\_FIB}(n-1) + \text{REC\_FIB}(n-2)\)
	\EndIf
\end{algo}

\vspace{5pt}

\noindent
L'equazione alle ricorrenze dell'algoritmo è la seguente:
\[T(n) = \begin{cases}
	0 & \text{se } n = 0,1 \\
	T(n-1) + T(n-2) + 1 & \text{se } n > 1
\end{cases}\]
Osservando che \(T(n-1) \geq T(n-2)\), si calcola un lower bound della complessità temporale che risulta essere doppiamente
esponenziale nella taglia dell'istanza, per cui fortemente inefficiente.
\begin{align*}
	T(n) \;\; &\geq \;\; 2T(n-2) + 1 \\[3pt]
	&= \;\; \Omega \! \left( 2^{n/2} \right) \; = \; \Omega \! \left( \sqrt{2}^n \right) \; = \; \Omega \! \left( 2^{2^{(\log_2 n) - 1}} \right) \; = \; \Omega \! \left( 2^{2^{\text{taglia(n)} - 1}} \right)
\end{align*}

\newpage

\subsubsection*{Soluzione iterativa con memoizzazione}
Siccome l'\(n\)-esimo numero di Fibonacci dipende solo dai due numeri precedenti, si osserva che è sufficiente tenere in
memoria i due numeri di Fibonacci più recenti appena calcolati per poter calcolare il successivo. Questo tipo di approccio
è anche detto bottom-up.

\begin{algo}{IT\_FIB}{n}{intero positivo $n \in \mathbb{Z}_{+,0}$}{$n$-esimo numero di Fibonacci $F(n) \in \mathbb{Z}_{+}$}
	\If{\(n = 0\) \textbf{or} \(n = 1\)}
		\State \Return \(n\)
	\EndIf

	\State
	\State \(F[0] \gets 0\)
	\State \(F[1] \gets 1\)
	\For{\(i \gets 2\) \textbf{to} \(n\)}
		\State \(F[i] \gets F[i-1] + F[i-2]\)
	\EndFor
	\State \Return \(F[n]\)
\end{algo}

~

\noindent
Siccome l'algoritmo non è ricorsivo, la complessità risulta più facile da calcolare e asintoticamente è esponenziale nella
taglia dell'istanza, per cui comunque poco efficiente.
\[T(n) = n + 1 = 2^{\log_2 n} - 1 = 2^\text{taglia(n)} - 1\]


\subsubsection*{Soluzione analitica con formula di Binet ed esponenziazione di un numero}
Riscrivendo la definizione ricorsiva di Fibonacci in forma matriciale e calcolando successivamente gli autovalori e autovettori
della matrice di trasferimento è possible convertire la definizione ricorsiva in una formula chiusa nota come formula di Binet.
\[\begin{pmatrix} F_n \\ F_{n-1} \end{pmatrix} = \begin{pmatrix} 1 & 1 \\ 1 & 0 \end{pmatrix} \begin{pmatrix} F_{n-1} \\ F_{n-2} \end{pmatrix} \qquad \rightarrow \qquad \begin{pmatrix} F_n \\ F_{n-1} \end{pmatrix} = \begin{pmatrix} 1 & 1 \\ 1 & 0 \end{pmatrix}^{n-1} \begin{pmatrix} F_{1} \\ F_{0} \end{pmatrix} = \begin{pmatrix} 1 & 1 \\ 1 & 0 \end{pmatrix}^{n-1} \begin{pmatrix} 1 \\ 0 \end{pmatrix}\]
\vspace{-7pt}
\[\text{autovalori ed autovettori} : \quad \lambda_1 = \frac{1 + \sqrt{5}}{2}, \quad \lambda_2 = \frac{1 - \sqrt{5}}{2}, \quad v_1 = \left(\frac{1 + \sqrt{5}}{2}, 1\right), \quad v_2 = \left(\frac{1 - \sqrt{5}}{2}, 1\right)\]
\vspace{-14pt}
\begin{align*}
	\begin{pmatrix} 1 & 1 \\ 1 & 0 \end{pmatrix}^{n-1} &= \begin{pmatrix} {v_1}^T \\ {v_2}^T \end{pmatrix} \cdot \begin{pmatrix} \lambda_1 & 0 \\ 0 & \lambda_2 \end{pmatrix}^{n-1} \cdot \begin{pmatrix} {v_1}^T & {v_2}^T \end{pmatrix}^{-1} = \\
	&= \begin{pmatrix} \frac{1 + \sqrt{5}}{2} & \frac{1 - \sqrt{5}}{2} \\[8pt] 1 & 1 \end{pmatrix} \cdot \begin{pmatrix} \left(\frac{1 + \sqrt{5}}{2}\right)^{n-1} & 0 \\[-2pt] 0 & \left(\frac{1 - \sqrt{5}}{2}\right)^{n-1} \end{pmatrix} \cdot \begin{pmatrix} \frac{1}{\sqrt{5}} & -\frac{1 - \sqrt{5}}{2\sqrt{5}} \\[8pt] -\frac{1}{\sqrt{5}} & \frac{1 + \sqrt{5}}{2\sqrt{5}} \end{pmatrix} = \\
	&= \begin{pmatrix} \frac{1}{\sqrt{5}} \cdot \left[ \left(\frac{1 + \sqrt{5}}{2}\right)^n - \left(\frac{1 - \sqrt{5}}{2}\right)^n \right] & \dots \\[8pt] \frac{1}{\sqrt{5}} \cdot \left[ \left(\frac{1 + \sqrt{5}}{2}\right)^{n-1} - \left(\frac{1 - \sqrt{5}}{2}\right)^{n-1} \right] & \dots \end{pmatrix} = \begin{pmatrix} \frac{1}{\sqrt{5}} \cdot \left( {\lambda_1}^n - {\lambda_2}^n \right) & \dots \\[8pt] \frac{1}{\sqrt{5}} \cdot \left( {\lambda_1}^{n-1} - {\lambda_2}^{n-1} \right) & \dots \end{pmatrix}
\end{align*}
\[\begin{pmatrix} F_n \\ F_{n-1} \end{pmatrix} = \begin{pmatrix} 1 & 1 \\ 1 & 0 \end{pmatrix}^{n-1} \begin{pmatrix} 1 \\ 0 \end{pmatrix} = \begin{pmatrix} \frac{1}{\sqrt{5}} \cdot \left( {\lambda_1}^n - {\lambda_2}^n \right) & \dots \\[8pt] \frac{1}{\sqrt{5}} \cdot \left( {\lambda_1}^{n-1} - {\lambda_2}^{n-1} \right) & \dots \end{pmatrix} \cdot \begin{pmatrix} 1 \\ 0 \end{pmatrix} \quad \rightarrow \quad F(n) = \frac{1}{\sqrt{5}} \cdot \left( {\lambda_1}^n - {\lambda_2}^n \right)\]
~

\noindent
La complessità della formula analitica di Binet è dominata dal calcolo delle potenze di numeri reali, che risultano
essere lineari nella taglia dell'istanza dell'esponente per come visto nell'algoritmo efficiente di esponenziazione,
per cui \(T(n) = \Theta(\log_2 n)\).

Siccome si lavora con numeri reali, però, servono analisi più approfondite per valutare sia la complessità, siccome
la taglia di un numero reale dipende dalla rappresentazione in virgola mobile utilizzata, sia l'accuratezza dei risultati,
siccome le rappresentazioni in virgola mobile introducono errori di approssimazione che non possono essere trascurati.
Per questo motivo l'algoritmo basato sulla formula di Binet si accenna solo dal punto di vista teorico.

\subsubsection*{Soluzione ricorsiva con memoizzazione}
È possibile sviluppare un approccio ricorsivo con memoizzazione, in cui la funzione ricorsiva costruisce e memorizza le
soluzioni intermedie in un vettore come nell'algoritmo iterativo, piuttosto che costruirle da zero ogni volta e dimenticarle.
La complessità di questo approccio è uguale a quella dell'algoritmo iterativo, ovvero esponenziale nella taglia dell'istanza.

L'implementazione è composta da una funzione di inizializzazione INIT\_FIB(\(n\)) che costruisce e inizializza il vettore
dove salvare le soluzioni intermedie e da una funzione ricorsiva REC\_FIB\_MEM(\(n\)) che costruisce le soluzioni intermedie
in modo ricorsivo memorizzandole nel vettore.

\begin{algo}{INIT\_FIB}{n}{intero positivo $n \in \mathbb{Z}_{+,0}$}{$n$-esimo numero di Fibonacci $F(n) \in \mathbb{Z}_{+}$}
	\If{\(n = 0\) \textbf{or} \(n = 1\)}
		\State \Return \(n\)
	\EndIf

	\State
	\State \(F[0] \gets 0\)
	\State \(F[1] \gets 1\)
	\For{\(i \gets 2\) \textbf{to} \(n\)}
		\State \(F[i] \gets 0\)
	\EndFor
	\State \Return \text{REC\_FIB\_MEM}\(n\)
\end{algo}

\begin{algo}{REC\_FIB\_MEM}{n}{intero positivo $n \in \mathbb{Z}_{+,0}$}{$n$-esimo numero di Fibonacci $F(n) \in \mathbb{Z}_{+}$}
	\If{\(F[n] = 0\)}
		\State \(F[n] \gets \text{REC\_FIB\_MEM}(n-1) + \text{REC\_FIB\_MEM}(n-2)\)
	\EndIf
	\State \Return \(F[n]\)
\end{algo}

\newpage

\subsection{Longest Common Subsequence - LCS}
\subsubsection*{Definizione di sottosequenza e di sottostringa di una stringa}
La sottosequenza di una stringa \(X\) è una successione di simboli non necessariamente contigui che mantengono lo stesso
ordinamento con cui si trovano nella stringa \(X\). Una sottostringa è una sottosequenza particolare in cui è richiesto
che i simboli devono essere contigui.
\[\begin{array}{c c}
	\displaystyle \binom{n}{k} = \frac{n!}{k! \, (n-k)!} & \quad \text{\# di sottosequenze di lunghezza \(k\) di una stringa di lunghezza \(n\)} \\[13pt]
	\displaystyle \sum_{k=0}^{n} \binom{n}{k} = 2^n & \quad \text{\# di tutte le sottosequenze di una stringa di lunghezza \(n\)}
\end{array}\]

\noindent
Si nota che per ottenere una sottosequenza di lunghezza \(k\), si scelgono a caso \(k\) simboli dalla stringa \(X\) di lunghezza
\(n\) senza tenere conto dell'ordine. È equivalente a cercare un sottoinsieme di \(k\) elementi da un insieme di \(n\)
elementi, per questo motivo si utilizza il coefficiente binomiale. L'ordine dei simboli nella sottosequenza è imposto a
posteriori della scelta, in base all'ordine con cui si trovano nella stringa \(X\).

\subsubsection*{Formalizzazione del problema}
Date due stringhe \(X\) e \(Y\) rispettivamente di lunghezza \(m\) e \(n\), si vuole calcolare la sottosequenza più lunga
\(Z\) che sia comune ad entrambe le stringhe \(X\) e \(Y\). La taglia dell'istanza è data da \(m + n\).

\subsubsection*{Descrizione dell'algoritmo LCS(\(X\), \(Y\)) ricorsivo con memoizzazione}
\begin{itemize}
	\item \textbf{caso base}: \\
	se una delle due stringhe \(X\) o \(Y\) è vuota, la sottosequenza comune più lunga è anch'essa vuota
	\item \textbf{passo ricorsivo - caso 1}: \\
	se le due stringhe \(X = \; \langle X', a \rangle\) e \(Y = \; \langle Y', a \rangle\) terminano con lo stesso simbolo \(a\), allora la sottosequenza
	comune più lunga è data dalla sottosequenza comune più lunga tra \(X'\) e \(Y'\) (\(\text{LCS}(X', Y')\)) con l'aggiunta del
	simbolo \(a\) alla fine
	\item \textbf{passo ricorsivo - caso 2}: \\
	se le due stringhe \(X = \; \langle X', a \rangle\) e \(Y = \; \langle Y', b \rangle\) terminano con simboli diversi \(a\) e \(b\), allora la sottosequenza
	comune più lunga sarà data da quella di lunghezza massima tra la sottosequenza comune più lunga tra \(X'\) e \(Y\) (\(\text{LCS}(X', Y)\)) e
	quella tra \(X\) e \(Y'\) (\(\text{LCS}(X, Y')\))
\end{itemize}
\vspace{5pt}

\noindent
Per identificare i prefissi di \(X\) e \(Y\) da passare alle chiamate ricorsive, si utilizzano rispettivamente gli
indici \(i\) e \(j\) che puntano all'ultimo simbolo di \(X\) e \(Y\) da considerare. Di seguito sono riportate la proprietà
di sottostruttura e il calcolo della lunghezza \(l\) della LCS indicate come equazioni alle ricorrenze.
\[Z(X_{0:i}, Y_{0:j}) \; = \text{LCS}(i, j) \; = \; \begin{cases}
	\; \varepsilon & \text{caso base} - \text{per } i = 0 \vee j = 0 \\[3pt]
	\; \text{LCS}(i-1, j-1) \;\; \cup \; \langle X_i \rangle & \text{caso 1} - \text{se } X_i = Y_j \\[3pt]
	\; \max \{\; \text{LCS}(i-1, j), \; \text{LCS}(i, j-1)\; \} & \text{caso 2} - \text{se } X_i \neq Y_j
\end{cases}\]
\[\text{len}(Z(X_{0:i}, Y_{0:j})) \; = \; l(i, j) \; = \; \begin{cases}
	0 & \text{caso base} - \text{per } i = 0 \vee j = 0 \\[3pt]
	l(i-1, j-1) + 1 & \text{caso 1} - \text{se }X_i = Y_j \\[3pt]
	\max \{ l(i-1, j), l(i, j-1) \} & \text{caso 2} - \text{se } X_i \neq Y_j
\end{cases}\]

\vspace{5pt}

\noindent
L'algoritmo viene definito con approccio bottom-up, ovvero si parte da un caso base semplice \(i = j = 0\) e si
espandono le soluzioni intermedie facendo variare progressivamente un indice alla volta fino ad arrivare alla soluzione
finale \(i = m\) e \(j = n\). Ogni soluzione intermedia viene memorizzata in una tabella bidimensionale, come illustrato
successivamente.

\newpage

\subsubsection*{Tabella dei risultati (struttura dati di memoizzazione)}
Le soluzioni intermedie vengono memorizzate in una tabella bidimensionale di dimensione \((m+1) \times (n+1)\) con le
seguenti caratteristiche:
\begin{itemize}
	\item la cella \((i,j)\) fa riferimento alla chiamata \(LCS(i, j)\)
	\item il valore memorizzato nella cella \((i,j)\) è la lunghezza \(l(i,j)\) della \(LCS(i, j)\)
	\item la soluzione finale \(LCS(X, Y)\) è contenuta nella cella \((m,n)\) (in basso a destra) della tabella
	\item in ogni cella è riportata una freccia che indica la cella riferita alla chiamata ricorsiva effettuata per
	calcolare il valore della cella stessa, con le seguenti regole:
	\begin{itemize}[topsep=-2pt]
		\item se \(X_i = Y_j\) la freccia è diagonale e punta alla cella \((i\!-\!1, \; j\!-\!1)\)
		\item se \(X_i \neq Y_j\) la freccia è orizzontale o verticale e punta rispettivamente alla cella \((i\!-\!1, \; j)\)
		o alla cella \((i, \; j\!-\!1)\) a seconda di quale delle due chiamate ricorsive ha restituito il valore più grande
		(alcune volte è indipendente scegliere tra orizzontale o verticale se i valori sono uguali)
	\end{itemize}
	\item per l'approccio di costruzione bottom-up, il popolamento della tabella parte dal caso base \((0,0)\) (in alto
	a sinistra) e prosegue per righe (approccio row-major) oppure per colonne (approccio column-major), in base a che
	indice si decide di far variare per primo
\end{itemize}

\noindent
Di seguito ne è riportato un esempio per le stringhe \(X = \, \langle A, B, C, B, B, D \rangle\) e \(Y = \; \langle A, D, C, C, B, D \rangle\).
\begin{center}
	\renewcommand{\arraystretch}{1.3}
	\setlength{\tabcolsep}{12pt}
	\begin{tabular}{c|c|c|c|c|c|c|c|}
	\diagbox{x}{y} & \(\begin{array}{c} \!\!\!\!Y_0\!\!\!\! \\ \varepsilon \end{array}\) & \(\begin{array}{c} \!\!\!\!Y_1\!\!\!\! \\ A \end{array}\) & \(\begin{array}{c} \!\!\!\!Y_2\!\!\!\! \\ D \end{array}\) & \(\begin{array}{c} \!\!\!\!Y_3\!\!\!\! \\ C \end{array}\) & \(\begin{array}{c} \!\!\!\!Y_4\!\!\!\! \\ C \end{array}\) & \(\begin{array}{c} \!\!\!\!Y_5\!\!\!\! \\ B \end{array}\) & \(\begin{array}{c} \!\!\!\!Y_6\!\!\!\! \\ D \end{array}\) \\
	\hline
	\(X_0: ~ \varepsilon\) & 0 & 0 & 0 & 0 & 0 & 0 & 0 \\
	\hline
	\(X_1: ~ A\) & 0 & $\nwarrow$ 1 & $\leftarrow$ 1 & $\leftarrow$ 1 & $\leftarrow$ 1 & 1 & 1 \\
	\hline
	\(X_2: ~ B\) & 0 & 1 & 1 & 1 & $\uparrow$ 1 & 2 & 2 \\
	\hline
	\(X_3: ~ C\) & 0 & 1 & 1 & 2 & $\nwarrow$ 2 & 2 & 2 \\
	\hline
	\(X_4: ~ B\) & 0 & 1 & 1 & 2 & $\uparrow$ 2 & 3 & 3 \\
	\hline
	\(X_5: ~ B\) & 0 & 1 & 1 & 2 & 2 & $\nwarrow$ 3 & 3 \\
	\hline
	\(X_6: ~ D\) & 0 & 1 & 2 & 2 & 2 & 3 & $\nwarrow$ 4 \\
	\hline
	\end{tabular}
\end{center}

\subsubsection*{Ricostruzione della LCS a partire dalle frecce della tabella}
La ricostruzione della sottosequenza comune più lunga \(Z\) equivale ad una visita in post-order dell'ordinamento topologico
delle frecce contenute nella tabella \(\underline{b}\). La prima chiamata deve essere effettuata con gli indici \(i = m\)
e \(j = n\) che puntano all'ultima cella della tabella. La visita termina quando si raggiunge una cella con indice \(i = 0\)
o \(j = 0\) (caso base).

\begin{algo}{PRINT\_LCS}{$\underline{X}$, $\underline{b}$, $i$, $j$}{stringa $\underline{X}$ per recupero elementi, tabella $\underline{b}$ con le frecce e indici $i$, $j$}{nessun return, stampa la LCS($X$, $Y$) a standard output}
	\If{\(m = 0\) \textbf{or} \(n = 0\)} \Comment{invocazione sul caso base \(\rightarrow\) nessuna azione necessaria}
		\State \Return \(-\)
	\EndIf
	\If{\(b[i, j] = \text{``}\nwarrow\text{''}\)} \Comment{caso 1 - cella diagonale}
		\State \text{PRINT\_LCS}(\(\underline{X}, \underline{b}, i-1, j-1\))
		\State print(\(X_i\))
	\ElsIf{\(b[i, j] = \text{``}\uparrow\text{''}\)} \Comment{caso 2 - cella verticale}
		\State \text{PRINT\_LCS}(\(\underline{X}, \underline{b}, i-1, j\))
	\Else \(\quad (\text{ovvero per }b[i, j] = \text{``}\leftarrow\text{''}\)) \Comment{caso 2 - cella orizzontale}
		\State \text{PRINT\_LCS}(\(\underline{X}, \underline{b}, i, j-1\))
	\EndIf
\end{algo}

\newpage

\subsubsection*{Implementazione iterativa bottom-up}
\vspace{-7pt}
\begin{algo}{IT\_LCS}{$\underline{X}$, $\underline{Y}$}{stringhe $\underline{X}$ e $\underline{Y}$ rispettivamente di lunghezza $m$ e $n$}{$(l(m, n), \underline{b})$ lunghezza e tabella con le frecce per ricostruire la LCS($X$, $Y$)}
	\State \(m \gets \text{len}(X)\)
	\State \(n \gets \text{len}(Y)\)
	\For {\(i \gets 0\)}{\(m\)} \(l[i, 0] \gets 0\) \Comment{inizializzazione caso base della prima riga} \EndFor
	\For {\(j \gets 0\)}{\(n\)} \(l[0, j] \gets 0\) \Comment{inizializzazione caso base della prima colonna} \EndFor

	\State
	\For{\(i \gets 1\)}{\(m\)} \Comment{popolamento della tabella row-major}
		\For{\(j \gets 1\)}{\(n\)}
			\If{\(X_i = Y_j\)} \Comment{passo ricorsivo - caso 1}
				\State \(l[i, j] \gets l[i-1, j-1] + 1 \quad-\quad b[i, j] \gets \text{``}\nwarrow\text{''}\)
			\ElsIf{\(l[i-1, j] \geq l[i, j-1]\)} \Comment{passo ricorsivo - caso 2 - cella verticale}
				\State \(l[i, j] \gets l[i-1, j] \quad-\quad b[i, j] \gets \text{``}\uparrow\text{''}\)
			\ElsIf{\(l[i-1, j] < l[i, j-1]\)} \Comment{passo ricorsivo - caso 2 - cella orizzontale}
				\State \(l[i, j] \gets l[i, j-1] \quad-\quad b[i, j] \gets \text{``}\leftarrow\text{''}\)
				\State \(b[i, j] \gets \text{``}\leftarrow\text{''}\)
			\EndIf
		\EndFor
	\EndFor
	\State \Return (\(l[m][n], \;\; \underline{b}\))
\end{algo}

\subsubsection*{Implementazione ricorsiva top-down con memoizzazione}
\vspace{-7pt}
\begin{algo}{INIT\_LCS}{$\underline{X}$, $\underline{Y}$}{stringhe $\underline{X}$ e $\underline{Y}$ rispettivamente di lunghezza $m$ e $n$}{$(l(m, n), \underline{b})$ lunghezza e tabella con le frecce per ricostruire la LCS($X$, $Y$)}
	\State \(m \gets \text{len}(X)\)
	\State \(n \gets \text{len}(Y)\)
	\If{\(m = 0\) \textbf{or} \(n = 0\)} \Comment{invocazione sul caso base \(\rightarrow\) nessuna azione necessaria}
		\State \Return \(0\)
	\EndIf
	\For {\(i \gets 0\)}{\(m\)} \(l[i, 0] \gets 0\) \Comment{inizializzazione caso base della prima riga} \EndFor
	\For {\(j \gets 0\)}{\(n\)} \(l[0, j] \gets 0\) \Comment{inizializzazione caso base della prima colonna} \EndFor
	\For {\(i \gets 1\)}{\(m\)} \Comment{inizializzazione con valore sentinella \(-1\)}
		\For {\(j \gets 1\)}{\(n\)} \Comment{per identificare i casi non ancora gestiti}
			\State \(l[i,j] \gets -1\)
		\EndFor
	\EndFor
	\State \Return \text{REC\_LCS}\((\underline{X}, \underline{Y},m, n)\)
\end{algo}

\begin{algo}{REC\_LCS}{$\underline{X}$, $\underline{Y}$, $m$, $n$}{stringhe $\underline{X}$ e $\underline{Y}$ e rispettive lunghezze $m$ e $n$}{$(l(m, n), \underline{b})$ lunghezza e tabella con le frecce per ricostruire la LCS($X$, $Y$)}
	\If{\(l[m, n] = -1\)} \Comment{caso non ancora gestito}
		\If{\(X_i = Y_j\)} \Comment{passo ricorsivo - caso 1}
			\State \(l[i, j] \gets 1 + \text{REC\_LCS}(\underline{X}, \underline{Y}, i-1, j-1) \quad-\quad b[i, j] \gets \text{``}\nwarrow\text{''}\)
		\ElsIf{\(\text{REC\_LCS}(\underline{X}, \underline{Y}, i\!-\!1, j) \geq \text{REC\_LCS}(\underline{X}, \underline{Y}, i, j\!-\!1)\)} \Comment{caso 2 - cella verticale}
			\State \(l[i, j] \gets l[i-1, j] \quad-\quad b[i, j] \gets \text{``}\uparrow\text{''}\)
		\Else \(\quad (\text{ovvero per }l[i-1, j] < l[i, j-1])\) \Comment{caso 2 - cella orizzontale}
			\State \(l[i, j] \gets l[i, j-1] \quad-\quad b[i, j] \gets \text{``}\leftarrow\text{''}\)
		\EndIf
	\EndIf
	\State \Return (\(l[m][n], \;\; \underline{b}\))
\end{algo}

\newpage

\subsubsection*{Complessità temporale}
La complessità temporale dell'algoritmo LCS è \(O(m \cdot n)\) sia per l'implementazione iterativa che per quella ricorsiva
con memoizzazione, per come di seguito spiegato:
\begin{itemize}
	\item nella versione iterativa, si visitano tutte le \(n \cdot m\) celle della tabella e per ogni cella viene fatto
	un numero costante di operazioni
	\item nella versione ricorsiva con memoizzazione, ogni cella della tabella deve essere inizializzata per cui il
	ragionamento è lo stesso della versione iterativa
\end{itemize}
Tale complessità, seppur quadratica per \(m = n\), è comunque molto più efficiente rispetto alla complessità esponenziale
dell'algoritmo banale brute-force che esplora tutte le possibili sottosequenze di \(X\) e \(Y\) per identificare quella
in comune più lunga.

Con alcuni accorgimenti è possibile ridurre la complessità temporale fino a \(O(m \cdot n / \log_2 n)\).

\subsubsection*{NO-SETH e Fine Grained Complexity}
Il problema LCS è definito come NO-SETH, ovvero è un algoritmo che se dovesse essere risolto in tempo polinomiale subquadratico
\(O(n^{2-\varepsilon})\), allora sarebbe possibile risolvere in tempo polinomiale molti altri problemi che attualmente hanno
solo soluzioni con complessità esponenziale. Questo studio è svolto nell'ambito della ``Fine Grained Complexity'', ovvero
una branca dello studio della complessità che si occupa di identificare i valori precisi delle complessità di algoritmi per
studiare come piccoli miglioramenti su alcuni algoritmi possano portare al passaggio da esponenziale a polinomiale per molti
altri problemi.

\subsubsection*{Complessità spaziale}
La complessità spaziale dell'algoritmo LCS è \(O(m \cdot n)\), ovvero è proporzionale al numero di celle della tabella in
cui vengono salvati i risultati. Corrisponde a tutte le possibili combinazioni di prefissi di \(X\) e \(Y\) che possono
essere passati alle chiamate ricorsive.

Se si è interessati a conoscere solo la lunghezza della LCS, è possibile ridurre la complessità spaziale a \(O(2 \cdot \min\{m, n\})\)
memorizzando solo le due righe (o colonne) più recenti della tabella (quella precedente e quella che si sta popolando),
in quando le equazioni alle ricorrenze dipendono solo dai valori delle celle adiacenti (diagonale, orizzontale e verticale).

\subsubsection*{Osservazioni sulla struttura dati per la memoizzazione}
Per gestire la memoizzazione è richiesta una struttura dati con le seguenti operazioni:
\begin{itemize}
	\item ricerca tramite la sottoistanza (coppia di indici (\(i, j\))) che restituisce la soluzione della sottoistanza cercata
	o \(-1\) se tale sottoistanza non è ancora stata gestita
	\item inserimento delle soluzioni per le sottoistanze risolte
\end{itemize}
La struttura più semplice che garantisce tali operazioni con complessità \(O(1)\) è una matrice bidimensionale, ma ha lo
svantaggio che lo spazio occupato è fisso di \(O(m \cdot n)\) indipendentemente dal numero di sottoistanze effettivamente
gestite.

Un'alternativa più efficiente a livello spaziale sono il dizionario o mappa (eventualmente semplificati in quanto non è
richiesta la cancellazione) in quanto occupano lo spazio necessario per memorizzare solo le sottoistanze effettivamente
gestite. Lo svantaggio è che hanno complessità di ricerca e inserimento di \(O(\log_2 (m \cdot n))\). Risultano, quindi,
più efficienti solo se il numero di sottoistanze effettivamente gestite è minore di \((m \cdot n)/\log_2 (m \cdot n)\).

\newpage

\subsection{Matrix Chain Multiplication - MCM}
\subsubsection*{Introduzione al problema}
Dati tre vettori \(p\), \(q\) e \(r\) rispettivamente di dimensione \(n \times 1\) (colonna), \(1 \times n\) (riga) e
\(n \times 1\) (colonna), si vuole trovare una parentesizzazione del loro prodotto per risolverlo in maniera efficiente.
\[(p \cdot q) \cdot r = (\text{colonna} \cdot \text{riga}) \cdot \text{colonna} = \text{matrice} \cdot \text{colonna} \quad \rightarrow \quad T(n) = \underbrace{\; n^2 \;}_{p_\text{colonna} \; \cdot \; q_\text{riga}} + \underbrace{\; 2n-1 \;}_{\text{matrice} \; \cdot \; r_\text{colonna}}\]
\[p \cdot (q \cdot r) = \text{colonna} \cdot (\text{riga} \cdot \text{colonna}) = \text{colonna} \cdot \text{scalare} \quad \rightarrow \quad T(n) = \underbrace{\; n \;}_{p_\text{colonna} \; \cdot \; \text{scalare}} + \underbrace{\; 2n-1 \;}_{q_\text{riga} \; \cdot \; r_\text{colonna}}\]
Si osserva che parentesizzando in modo diverso, si ottengono complessità decisamente diverse. Elencare tutte le possibili
parentesizzazioni richiederebbe costi esponenziali, per cui è necessario identificare un approccio più efficiente.

\subsubsection*{Formalizzazione del problema}
L'input del problema è una sequenza di \(n+1\) interi \(\langle  p_0, p_1, \dots, p_{n-1}, p_n \rangle\) che rappresentano
le dimensioni delle \(n\) matrici da moltiplicare, secondo la relazione \(A_i \in \mathbb{C}^{p_{i-1} \times p_i}\)
con \(i = 1, \dots, n\).

Inoltre si rappresenta una parentesizzazione di una sequenza di matrici \(\langle  A_i, A_{i+1}, \dots, A_j \rangle\) attraverso
la coppia di indici \(i\) e \(j\) che identificano la prima e l'ultima matrice della sottostringa.

\subsubsection*{Idea risolutiva D\&C e presenza di istanze sovrapposte}
Data una stringa di \(j-i+1\) matrici \(\langle A_i, A_{i+1}, \dots, A_j \rangle\) delimitate dagli indici \(i\) e \(j\),
si identifica la posizione \(k\) ottimale in cui dividere o parentesizzare tale stringa generando due sottostringhe
\(\langle A_i, A_{i+1}, \dots, A_k \rangle\) e \(\langle A_{k+1}, A_{k+2}, \dots, A_j \rangle\) delimitate rispettivamente dalle
coppie di indici \((i, k)\) e \((k+1, j)\). Una volta fatto ciò, si calcola ricorsivamente il costo dei prodotti
necessari a risolvere le due sottostringhe e ottenere le due matrici \(A_{i:k}\) e \(A_{k+1:j}\) a cui poi si aggiungerà
il costo necessario a moltiplicarle tra loro per ottenere la matrice \(A_{i:j}\) corrispondente al prodotto di tutte
le matrici della stringa originale.

Si definisce la funzione \(c[i, j]\) come il costo minimo della moltiplicazione delle matrici della sottostringa
\(\langle A_i, A_{i+1}, \dots, A_j \rangle\) supponendo di aver scelto il valore \(k\) ottimo per minimizzare tale costo.
È utile notare che il costo del prodotto tra le due matrici \(A_{i:k} \in \mathbb{C}^{p_{i-1} \times p_k}\) e
\(B_{k+1:j} \in \mathbb{C}^{p_k \times p_j}\) è dato da \(p_{i-1} \cdot p_k \cdot p_j\).

\[c[i, j] = \begin{cases}
	\; 0 & \text{se } i = j \\
	\; \displaystyle \min_{i \leq k < j} \left\{ \; c[i, k] \; + \; c[k+1, j] \; + \; p_{i-1} \cdot p_k \cdot p_j \; \right\} & \text{se } i < j
\end{cases} \]

\noindent
Si osserva, però, che alcune sottostringhe di matrici possono comparire più volte nell'albero delle chiamate ricorsive,
per cui un approccio puramente ricorsivo con D\&C risulterebbe inefficiente. Si preferisce quindi sfruttare il paradigma
della programmazione dinamica.

\subsubsection*{Soluzione tramite memoizzazione}
Come per il problema LCS, si memorizzano le soluzioni delle sottoistanze identificate dalle coppie di indici \((i, j)\)
all'interno delle celle di una tabella bidimensionale di dimensione \(n \times n\). Si osserva che la tabella ha
le seguenti caratteristiche:
\begin{itemize}
	\item siccome \(i \leq j\), le uniche celle della tabella che vengono effettivamente utilizzate sono quelle sopra
	la diagonale principale, per cui la tabella è triangolare superiore
	\item le celle lungo la diagonale principale \(c[i, j]\) con \(i=j\) contengono il costo \(0\), ovvero il caso base
	\item le celle appena sopra alla diagonale principale \(c[i, j]\) con \(j-i = 1\) riflettono tutte le sottostringhe
	di due matrici, le celle appena sopra a queste riflettono tutte le sottostringhe di tre matrici e così via
	\item la cella \(c[1, n]\) in alto a destra contiene la soluzione finale, ovvero il costo minimo del prodotto di tutte
	le matrici della stringa in input
\end{itemize}

\subsubsection*{Implementazione iterativa}
\vspace{-7pt}
\begin{algo}{IT\_MCM}{$\underline{p}$}{vettore $\underline{p}$ di dimensione $n+1$ con le dimensioni delle matrici}{$m[1, n]$ costo minimo come definito da $c[1,n]$, matrice $\underline{s}$ con i valori di $k$ ottimali}
	\State \(n \gets \text{len}(p) - 1\)
	\For{\(i \gets 1\)}{\(n\)} \(m[i, i] \gets 0\) \Comment{inizializzazione caso base della diagonale principale} \EndFor
	\For{\(l \gets 2\)}{\(n\)} \Comment{per ogni lunghezza \(l\) della sottostringa di matrici}
		\For{\(i \gets 1\)}{\(n-l+1\)} \Comment{discesa lungo la diagonale delle sottostringhe di lunghezza \(l\)}
			\State \(j \gets i + l - 1\)
			\State \(m[i, j] \gets +\infty\) \Comment{inizializzazione con valore sentinella}
			\For{\(k \gets i\)}{\(j-1\)} \Comment{per ogni possibile punto di divisione \(k\)}
				\State \(q \gets m[i, k] + m[k+1, j] + p_{i-1} \cdot p_k \cdot p_j\) \Comment{calcolo del costo per la divisione in \(k\)}
				\If{\(q < m[i, j]\)} \Comment{se il costo calcolato è minore del costo attuale memorizzato}
					\State \(m[i, j] \gets q\) \Comment{aggiornamento del costo minimo per la sottostringa \(i, j\)}
					\State \(s[i, j] \gets k\) \Comment{memorizzazione del punto di divisione ottimo per la sottostringa \(i, j\)}
				\EndIf
			\EndFor
		\EndFor
	\EndFor
	\State \Return \((m[1, n], \underline{s})\)
\end{algo}

\subsubsection*{Alcune osservazioni sull'implementazione iterativa}
\begin{itemize}
	\item il popolamento della tabella avviene per diagonali, partendo da quella principale e procedendo verso l'alto a destra:
	il primo ``for'' varia la lunghezza \(l\) (sceglie la diagonale), il secondo ``for'' varia l'indice \(i\) (sceglie la riga
	della cella da popolare) e l'indice \(j\) viene calcolato di conseguenza in base a \(i\) e \(l\)
	\item in presenza di parentesizzazioni equivalenti (ad esempio per matrici quadrate), l'algoritmo individua quella con
	valore di \(k\) minimo, in quanto utilizzando l'operatore \(<\) e non \(\leq\) nell'if a riga 9, le soluzioni equivalenti
	non vengono considerate migliorative dell'ottimo localmente memorizzato
	\item al posto di avere due tabelle triangolari superiori \(m\) e \(s\), è possibile utilizzare la stessa tabella \(m\)
	popolata superiormente con i costi e inferiormente con i valori di \(k\) ottimali
\end{itemize}

\subsubsection*{Complessità spaziale e temporale}
La complessità temporale dell'algoritmo MCM è \(O(n^3)\) in quanto si visitano circa la metà delle celle della tabella che
asintoticamente sono \(O(n^2)\) e per ciascuna di esse si esegue un ciclo su \(k\) che varia da \(i\) a \(j\) con
\(i, j \in [1, n]\) in cui si svolgono operazioni in tempo costante.

La complessità spaziale dell'algoritmo MCM è \(O(n^2)\) in quanto le tabelle impiegate devono avere dimensione \(n \times n\)
per memorizzare i costi e le posizioni di parentesizzazione ottimali per tutte le possibili coppie di indici \(i\) e \(j\).
A differenza dell'algoritmo di LCS, il valore di una cella \(m[i, j]\) non dipende esclusivamente dalle celle adiacenti, per
cui non si possono fare ottimizzazioni a livello spaziale.

\newpage

\subsection{Longest Palindrome Subsequence - LPS}
\subsubsection*{Definizione di sequenza palindroma}
Una sequenza \(Z = \; \langle z_1, z_2, \dots, z_n \rangle\) si dice palindroma se \(z_{1+h} = z_{n-h}\) per ogni \(h = 0, 1, \dots, n\!-\!1\).

\noindent
Si ricorda che una sottosequenza non necessariamente deve contenere simboli contigui.

\subsubsection*{Definizione della proprietà di sottostruttura}
Per definire la più lunga sottosequenza palindroma di una sequenza \(Z_{i:j}\), si sceglie di confrontare i due simboli
agli estremi \(z_i\) e \(z_j\):
\begin{itemize}
	\item se sono uguali, allora la sottosequenza palindroma conterrà tali due simboli oltre alla più lunga sottosequenza
	palindroma della sottostringa \(Z_{i+1:j-1}\)
	\item se sono diversi, allora la sottosequenza palindroma non conterrà almeno uno dei due simboli, per cui si sceglierà
	la più lunga tra quella della sottostringa \(Z_{i+1:j}\) e quella della sottostringa \(Z_{i:j-1}\)
\end{itemize}
Di seguito si formalizza la proprietà di sottostruttura in termini di equazioni alle ricorrenze.
\[LPS(Z_{i:j}) = \begin{cases}
	\; \left\{ z_i \right\} & \text{se } i = j \\[3pt]
	\; \left\{ z_i, z_j \right\} & \text{se } j - i = 1 \text{ e } z_i = z_j \\[3pt]
	\; \left\{ z_i \right\} \text{ oppure } \left\{ z_j \right\} & \text{se } j - i = 1 \text{ e } z_i \neq z_j \\[3pt]
	\; \left\{ z_i \right\} \;\; \cup \;\; LPS(Z_{i+1:j-1}) \;\; \cup \;\; \left\{ z_j \right\} & \text{se } j-i > 1 \text{ e } z_i = z_j \\[3pt]
	\; \displaystyle \argmax_{\substack{LPS(Z_{i+1:j}) \\ LPS(Z_{i:j-1})}} \left\{ \; \text{len}(LPS(Z_{i+1:j})), \; \text{len}(LPS(Z_{i:j-1})) \; \right\} & \text{se } j-i > 1 \text{ e } z_i \neq z_j
\end{cases}\]
In alternativa, si può definire la proprietà di sottostruttura in termini di lunghezza della sottosequenza palindroma,
nel caso in cui si è interessati solo alla sua lunghezza e non ai simboli che la compongono.
\[p(i, j) = \begin{cases}
	\; 1 & \text{se } i = j \\
	\; 2 & \text{se } j - i = 1 \text{ e } z_i = z_j \\
	\; 1 & \text{se } j - i = 1 \text{ e } z_i \neq z_j \\
	\; p(i+1, j-1) + 2 & \text{se } j-i > 1 \text{ e } z_i = z_j \\
	\; \displaystyle \max \left\{ p(i+1, j), p(i, j-1) \right\} & \text{se } j-i > 1 \text{ e } z_i \neq z_j
\end{cases}\]

\subsubsection*{Note sull'implementazione}
L'implementazione iterativa bottom-up a pagina seguente è molto simile a quella dell'algoritmo per la Longest Common
Subsequence (LCS) tra due stringhe. Come per l'LCS, si utilizza una tabella \(p\) per memorizzare la lunghezza della
sottosequenza palindroma più lunga per ogni sottostringa \(Z_{i:j}\) e una tabella \(b\) per memorizzare a che cella
adiacente è associata alla sottoistanza generata. La principale differenza è che la tabella è riempita a diagonali
come nella Matrix Chain Multiplication (MCM).

Una volta completato l'algoritmo, per ricostruire la sottosequenza palindroma più lunga, si parte dalla cella \(p[1, n]\)
che contiene la lunghezza della sottosequenza palindroma più lunga per l'intera stringa \(Z\) e si procede scendendo
nella tabella \(p\) seguendo le direzioni memorizzate nella tabella \(b\) fino ad arrivare al caso base nelle celle
della diagonale \(p[i, i]\).

\subsubsection*{Complessità computazionale e spaziale}
La complessità computazionale è \(T(n) = \Theta(n^2)\) perché si devono visitare tutte le celle (o al più metà) della
tabella \(p\) di dimensione \(n \times n\) e per ogni cella si effettuano un numero costante di operazioni.

La complessità spaziale è \(S(n) = \Theta(n^2)\) per la tabella \(p\) di dimensione \(n \times n\) e per la tabella \(b\)
di dimensioni analoghe. Come per la LCS, se si è interessati solo alla lunghezza della sottosequenza palindroma più lunga
è possibile ridurre la complessità spaziale a \(S(n) = \Theta(n)\) memorizzando solo le due diagonali della tabella \(p\)
necessarie per soddisfare le dipendenze del calcolo iterativo.

\newpage
\subsubsection*{Implementazione iterativa}
\vspace{-7pt}
\begin{algo}{IT\_LPS}{$\underline{x}$}{sequenza $\underline{x}$}{lunghezza della LPS di $\underline{x}$ e tabella delle transizioni $\underline{b}$ per ricostruire la LPS}
	\State \(n \gets \text{len}(\underline{x})\)
	\For{\(i \gets 1\)}{\(n-1\)} \Comment{inizializzazione dei casi base per le sottosequenze di lunghezza 1 e 2}
		\State \(p[i, i] \gets 1\)
		\If{\(x[i] = x[i+1]\)} \(p[i, i+1] \gets 2\)
		\Else ~ \(p[i, i+1] \gets 1\) \EndIf
	\EndFor
	\State \(p[n,n] \gets 1\)
	\State
	\For{\(l \gets 3\)}{\(n\)} \Comment{calcolo iterativo delle sottosequenze di lunghezza crescente}
		\For{\(i \gets 1\)}{\(n-l+1\)}
			\State \(j \gets i + l - 1\)
			\If{\(x[i] = x[j]\)}
				\State \(p[i, j] \gets p[i+1, j-1] + 2\)
				\State \(b[i, j] \gets \text{``} \swarrow \text{''}\)
			\ElsIf{\(p[i+1, j] \geq p[i, j-1]\)}
				\State \(p[i, j] \gets p[i+1, j]\)
				\State \(b[i, j] \gets \text{``} \downarrow \text{''}\)
			\Else \(\quad (\text{ovvero per } p[i+1, j] < p[i, j-1])\)
				\State \(p[i, j] \gets p[i, j-1]\)
				\State \(b[i, j] \gets \text{``} \leftarrow \text{''}\)
			\EndIf
		\EndFor
	\EndFor
	\State \Return \(p[1, n], \;\; \underline{b}\)
\end{algo}

\begin{algo}{PRINT\_LPS}{$\underline{x}$, $\underline{b}$, $i$, $j$}{sequenza $\underline{x}$, tabella delle transizioni $\underline{b}$, indici $i$ e $j$}{nessun return, stampa la LPS($\underline{x}$) a standard output}
	\If{\(i = j\)} print(\(x[i]\)) \Comment{caso base: sottosequenza di lunghezza 1} \EndIf
	\If{\(j - i = 1\)} \Comment{caso base: sottosequenza di lunghezza 2}
		\If{\(x[i] = x[j]\)} print(\(x[i], x[j]\))
		\Else ~ print(\(x[i]\)) \EndIf
	\EndIf
	\State
	\If{\(b[i, j] = \text{``}\swarrow\text{''}\)} \Comment{simboli agli estremi uguali}
		\State print(\(x[i]\))
		\State \text{PRINT\_LPS}(\(\underline{x}, \underline{b}, i+1, j-1\))
		\State print(\(x[j]\))
	\ElsIf{\(b[i, j] = \text{``}\downarrow\text{''}\)} \Comment{caso 1 - cella sotto}
		\State \text{PRINT\_LPS}(\(\underline{x}, \underline{b}, i+1, j\))
	\Else \(\quad (\text{ovvero per }b[i, j] = \text{``}\leftarrow\text{''}\)) \Comment{caso 2 - cella a sinistra}
		\State \text{PRINT\_LPS}(\(\underline{x}, \underline{b}, i, j-1\))
	\EndIf
\end{algo}

\section{Paradigma Greedy}
\subsection{Proprietà della scelta greedy}
\subsubsection*{Scelta greedy}
Il paradigma greedy è un caso particolare di divide and conquer, in cui la soluzione di un'istanza viene costruita ad ogni
passo induttivo, attraverso la scelta di un elemento ottimo locale, con l'obiettivo di ottenere una soluzione globale ottima.
La scelta effettuata è detta scelta greedy e spesso è estremamente banale (con costo costante), la parte complessa è
dimostrare che la scelta greedy porta ad una soluzione globale ottima.

\subsubsection*{Proprietà di scelta greedy}
La proprietà di scelta greedy è la condizione che deve soddisfare una scelta greedy per garantire che la scelta effettuata
localmente non comprometta l'ottimalità della soluzione globale. Una scelta greedy riduce sia l'istanza da risolvere che
lo spazio generale delle soluzioni (ottimali e subottimali). Per portare alla soluzione (e soddisfare la proprietà di scelta
greedy) è necessario che lo spazio delle soluzioni, dopo la riduzione, contenga ancora almeno una soluzione ottima globale. 

\subsubsection*{Dimostrazione della correttezza di un algoritmo greedy}
A livello pratico, per dimostrare che una scelta greedy soddisfa la proprietà di scelta greedy, è necessario dimostrare
che esiste una soluzione ottima globale che soddisfi la scelta greedy. Spesso si considera una generica soluzione
ottima globale e, se tale soluzione non contiene la scelta greedy, di dimostra che è possibile modificarla senza
comprometterne l'ottimalità, in modo che soddisfi la scelta greedy.

\subsubsection*{Algoritmi greedy subottimali}
Esistono alcuni problemi per cui sono stati sviluppati algoritmi greedy subottimali, che non garantiscono in assoluto il
raggiungimento della soluzione ottima globale, ma la raggiungono o ci si avvicinano molto con alta probabilità. Questi
algoritmi risultano essere molto più veloci degli algoritmi ottimi e sono impiegati in cui si è disposti a sacrificare
un po' di ottimalità in cambio di una maggiore velocità di esecuzione.

\subsection{Scheduling di attività}
\subsubsection*{Formalizzazione del problema}
Data una sequenza \(S=\left\{ a_1, a_2, \dots a_n \right\}\) di \(n\) attività \(a_i : [s_i, f_i)\), identificate dal proprio
istante di inizio \(s_i\) e di fine \(f_i\) con \(s_i < f_i\), si vuole individuare un sottoinsieme fattibile
\(A \subseteq S\) di dimensione massima.

\begin{itemize}[topsep=5pt]
	\item un sottoinsieme \(A \subseteq S\) di attività è fattibile se composto da attività mutualmente compatibili
	\item due attività \(a_i\) e \(a_j\) sono compatibili se e solo se i loro intervalli di tempo non si sovrappongono:
	\vspace{-5pt}
	\[a_i : [s_i, f_i) \cap a_j : [s_j, f_j) = \emptyset \qquad \Longleftrightarrow \qquad s_i \geq f_j \;\; \vee \;\; s_j \geq f_i\]
\end{itemize}

\subsubsection*{Definizione della scelta greedy}
Data una sequenza \(S\) di \(n\) attività, si costruisce l'insieme soluzione \(A\) scegliendo l'attività \(a_i \in S\) che
termina per prima (con \(f_i = \min_j f_j\)) unita all'insieme soluzione \(A_\text{sub}\) della sottoistanza \(S_\text{sub}\)
ottenuta partendo da \(S\) e rimuovendo tutte le attività incompatibili con l'attività scelta \(a_i\).

Di seguito si formalizza la scelta greedy attraverso la proprietà di sottostruttura espressa in termini di equazioni alle
ricorrenze.
\[A(S) = \begin{cases}
	\left\{ a_1 \right\} & \text{se } |S| = n = 1 \\[5pt]
	\displaystyle \left\{ \left\{ a_k \in S \text{ con } k = \argmin f_i \right\} \;\; \cup \;\; A( \, S - \left\{ a_l \in S \text{ incompatibili con } a_k \right\} \,) \right\} & \text{se } |S| = n > 1
\end{cases}\]

\subsubsection*{Algoritmo ricorsivo ottimizzato basato sulla scelta greedy}
Per rendere più pratica l'operazione di scelta dell'attività ottima, si ordina la sequenza \(S\) in ordine crescente rispetto
agli istanti di fine \(f_i\) delle varie attività, con costo di ordinamento \(\Theta(n \log_2 n)\). In questo modo l'attività
da scegliere sarà sempre la prima della sequenza \(S\) ordinata.

Inoltre si osserva che non è necessario visitare tutta la sequenza \(S\) per rimuovere le tutte le attività incompatibili
ogni volta che viene scelta un'attività \(a_i\). È sufficiente fermarsi non appena si trova la prima attività compatibile
\(a_j\) che sarà l'attività successiva da aggiungere all'insieme soluzione \(A\). Le attività incompatibili con \(a_i\)
saranno anche incompatibili con \(a_j\), per cui verranno rimosse successivamente senza influenzare la scelta ottima.

\begin{algo}{REC\_GREEDY\_SEL}{$\underline{s}$, $\underline{f}$, i}{$\underline{s}$ e $\underline{f}$ istanti di inizio e fine ordinati per $f_i$, indice $i$ della prima attività compatibile}{$A = \left\{ a_{i_1}, a_{i_2}, \dots a_{i_k} \right\} \subseteq S$ fattibile e di dimensione massima}
	\State \(n \gets \text{len}(\underline{s})\)
	\If{\(i > n\)} \Return \(\varnothing\) \Comment{caso base: non ci sono più attività da considerare} \EndIf
	\State \(A \gets \left\{ i \right\}\) \Comment{scelta greedy: si sceglie la prima attività della sequenza}
	\State \(j \gets i + 1\)
	\While{\(j \leq n\) \textbf{and} \(\underline{s}[j] < \underline{f}[i]\)} \(j \gets j + 1\) \Comment{si saltano le attività incompatibili con \(a_i\)} \EndWhile
	\State \Return \(A \cup \text{REC\_GREEDY\_SEL}(\underline{s}, \underline{f}, j)\) \Comment{proprietà di sottostruttura con chiamata ricorsiva}
\end{algo}

\subsubsection*{Dimostrazione della correttezza - verifica della proprietà greedy}
Di seguito si dimostra che la scelta greedy ideata rispetta la proprietà greedy, ovvero che esiste una soluzione ottima
\(A^*\) che rispetta la scelta greedy dell'algoritmo.

Si suppone di avere una generica soluzione ottima \(A*\):
\begin{itemize}
	\item se la prima attività \(a_1\) della soluzione ottima \(A^*\) è \(a_k\) con \(k = \argmin_i f_i\), allora la scelta
	greedy è banalmente rispettata
	\item se la prima attività \(a_1\) della soluzione ottima \(A^*\) non è \(a_k\) con \(k = \argmin_i f_i\), vuol dire che
	\(f_1 > f_k\), per cui è possibile sostituire \(a_1\) con \(a_k\) senza perdere fattibilità, ottenendo una nuova soluzione
	che rispetta la scelta greedy e rimane soluzione ottima
\end{itemize}

\subsubsection*{Analisi della complessità}
Escludendo il costo dell'ordinamento iniziale \(\Theta(n \log_2 n)\), la complessità dell'algoritmo è data soltanto dalle
operazioni di confronto durante la visita della sequenza \(S\) per saltare le attività incompatibili. Siccome le attività
della sequenza vengono visitate in maniera sequenziale e ad ogni visita vengono effettuate un numero costante di operazioni,
la complessità asintotica dell'algoritmo è \(T(n) = \Theta(n)\).

\newpage

\subsection{Codice di Huffman}
\subsubsection*{Introduzione al codice di Huffman}
Il codice di Huffman è un codice a prefisso utilizzato in teoria dell'informazione per codificare simboli con una certa
distribuzione di probabilità nota, in sequenze di bit dette parole di codice di lunghezza variabile. Ha la caratteristica
di essere ottimo, ovvero di minimizzare la lunghezza media dei codici generati.

Essendo un codice a prefisso, il codice di Huffman garantisce che nessuna parola di codice associata ad un simbolo sia
prefisso di un'altra parola associata ad un altro simbolo. Questo aspetto aiuta in fase di decodifica perché, appena si
individua una parola di codice associata ad un simbolo, la si può decodificare immediatamente e senza ambiguità.

Le prestazioni di un codice di Huffman sono misurate in termini di lunghezza media dei codici generati ed è definita come
la lunghezza media delle parole di codice, pesata in base alla loro probabilità:
\[L = \sum_{i=0}^{\# \; \text{parole}} p_i \cdot l_i \qquad \text{con} \qquad \begin{aligned}
	p_i & = \text{probabilità di occorrenza del simbolo } i \\[3pt]
	l_i & = \text{lunghezza della parola di codice associata al simbolo } i
\end{aligned}\]

\subsubsection*{Rappresentazione tramite Trie}
Il trie è un albero binario pieno (ogni nodo interno ha due figli), non completo (le foglie possono essere a livelli diversi)
utilizzato per rappresentare le parole di codice generate dal codice di Huffman. Ha le seguenti caratteristiche:
\begin{itemize}
	\item ogni nodo interno del trie ha esattamente due figli, quello di sinistra raggiungibile tramite un arco etichettato
	con un bit 0 e quello di destra raggiungibile tramite un arco etichettato con un bit 1
	\item ogni foglia del trie rappresenta un simbolo, la cui parola di codice è data dalla concatenazione dei bit etichettati
	sugli archi del percorso che va dalla radice alla foglia
\end{itemize}
Si osservano le seguenti proprietà del trie:
\begin{itemize}
	\item se una parola di codice fosse prefisso di un'altra, allora la parola più corta sarebbe rappresentata da un
	nodo interno del trie lungo il percorso che porta dalla radice alla foglia che rappresenta la parola più lunga; ciò è
	correttamente impossibile per come definito il trie
	\item si possono scambiare le due etichette \(0,1\) degli archi che collegano i due figli al padre senza alterare
	le prestazioni del codice
	\item si possono scambiare simboli che si trovano allo stesso livello del trie senza alterare le prestazioni del codice,
	anche se hanno nodi padre diversi
\end{itemize}

\subsubsection*{Algoritmo greedy e costruzione dell'albero}
La costruzione dell'albero avviene in maniera bottom-up:
\begin{itemize}
	\item si parte inizializzando una foresta di \(n\) alberi, ognuno costituito da un singolo nodo foglia che rappresenta
	un simbolo con la relativa probabilità di occorrenza
	\item ad ogni passo si scelgono i due alberi con le probabilità della radice più piccole, e si crea un nuovo nodo padre
	comune alle due radici, con probabilità pari alla somma delle probabilità delle due radici; per convenzione si collega
	l'albero con probabilità più piccola alla sinistra assegnando un bit 0 all'arco, e l'albero con probabilità più grande
	a destra assegnando un bit 1 all'arco
	\item si ripete il passo precedente fino a quando tutte le foreste non si sono unite in un unico albero, che rappresenta
	il trie del codice di Huffman e la radice ha probabilità pari a 1
\end{itemize}

La scelta greedy è quella di prendere i due alberi le cui radici hanno le probabilità più piccole.

\newpage

\subsubsection*{Implementazione greedy iterativa con priority queue}
\vspace{-7pt}
\begin{algo}{IT\_MCM}{$\underline{c}$, $\underline{f}$}{vettori $\underline{c}$ e $\underline{f}$ rispettivamente degli $n$ simboli e delle relative frequenze}{riferimento alla radice del trie che rappresenta il codice di Huffman}
	\State \(n \gets \text{len}(\underline{c})\)
	\State \(Q \gets \text{PriorityQueue}()\) \Comment{creazione di una coda di priorità vuota}
	\ForAll{\(c\)}{\(\underline{c}\)} \Comment{inizializzazione della coda di priorità con i nodi foglia}
		\State \(z \gets \text{new\_node}(\textit{symbol}: c, \;\; \textit{freq}: \underline{f}[c], \;\; \textit{left\_node}: \text{null}, \;\; \textit{right\_node}: \text{null})\)
		\State \(Q.\text{insert}(z.\text{freq}, z)\)
	\EndFor
	\For{\(i \gets 1\)}{\(n-1\)} \Comment{ad ogni iterazione diminuisco di 1 il numero di alberi nella foresta}
		\State \(x \gets Q.\text{extract\_min}() \;\; - \;\; y \gets Q.\text{extract\_min}()\)
		\State \(z \gets \text{new\_node}(\textit{symbol}: \text{null}, \;\; \textit{freq}: x.\text{freq} + y.\text{freq}, \;\; \textit{left\_node}: x, \;\; \textit{right\_node}: y)\)
		\State \(Q.\text{insert}(z.\text{freq}, z)\)
	\EndFor
	\State \Return \(Q.\text{extract\_min}()\)
\end{algo}
\vspace{-18pt}

\subsubsection*{Analisi della complessità computazionale}
Si sceglie di usare un heap binario per implementare la coda di priorità, che garantisce un costo \(\Theta(\log_2 n)\) per ogni
inserimento ed estrazione. La complessità dell'algoritmo è data da:
\begin{itemize}
	\item \(n + n-1 = 2n-1\) inserimenti nella coda di priorità, con costo \(\Theta(\log_2 n)\) per ogni inserimento
	\item \(2(n-1) + 1 = 2n-1\) estrazioni dalla coda di priorità, con costo \(\Theta(\log_2 n)\) per ogni estrazione
	\item costo lineare \(\Theta(n)\) dei due for, al netto delle operazioni di inserimento ed estrazione
\end{itemize}
\vspace{-7pt}

\[T(n) = (2n-1) \cdot T_{\, \text{insert()}} + (2n-1) \cdot T_{\, \text{extract\_min}()} + \Theta(n) = \Theta(n \log_2 n)\]

\noindent
Si osserva che, i primi \(n\) inserimenti delle foglie possono essere fatti in tempo \(\Theta(n)\) tramite l'algoritmo di
costruzione di un heap dato un array di \(n\) elementi. Siccome, però, l'algoritmo di Huffman richiede almeno \(n\) estrazioni
e non si dispone di algoritmi migliori di \(\Theta(\log_2 n)\) per l'estrazione, la complessità dell'algoritmo rimane \(\Theta(n \log_2 n)\).

\subsubsection*{Dimostrazione della correttezza}
Di seguito si dimostra che la scelta greedy di estrarre i due nodi di frequenza minima \(a\) e \(b\) per primi per
posizionali al livello più basso dell'albero, porta ad una soluzione ottima globale.
\begin{itemize}
	\item si considera un albero ottimo \(T\) con costo ottimo \(B(T) = \sum_{i \in T} f(i) \cdot d(i) \leq B(T^*)\) minore
	di qualsiasi altro albero \(T^*\) in cui:
	\begin{itemize}[topsep=-2pt]
		\item due foglie \(a\) e \(b\) che si trovano alla massima profondità dell'albero, uniti da un nodo padre \(p\)
		\item due foglie \(x\) e \(y\) con frequenze minime \(\max\left\{ f(x), f(y) \right\} \leq \min \left\{ f(a), f(b) \right\}\)
	\end{itemize}
	\item se \(a \equiv x\) e \(b \equiv y\), la scelta greedy è banalmente rispettata
	\item altrimenti si ipotizza di scambiare \(a\) con \(x\) e \(b\) con \(y\), ottenendo un nuovo albero \(T'\) e provocando
	una variazione del costo dell'albero \(\Delta B = B(T') - B(T)\) pari a:
	\begin{align*}
		\Delta B &= (f(a) \, d(x) \!+\! f(b) \, d(y) \!+\! f(x) \, d(a) \!+\! f(y) \, d(b)) - (f(a) \, d(a) \!+\! f(b) \, d(b) \!+\! f(x) \, d(x) \!+\! f(y) \, d(y)) \\
		&= f(a) \, (d(x) \!-\! d(a)) + f(y) \, (d(b) \!-\! d(y)) + f(x) \, (d(a) \!-\! d(x)) + f(b) \, (d(y) \!-\! d(b)) \\
		&= \underbrace{(f(a) - f(x))}_{\geq 0} \, \underbrace{(d(x) - d(a))}_{\leq 0} + \underbrace{(f(b) - f(y))}_{\geq 0} \, \underbrace{(d(y) - d(b))}_{\leq 0} \leq 0 \quad \Rightarrow \quad B(T') \leq B(T)
	\end{align*}
	\item da \(B(T) \leq B(T')\) per l'ipotesi iniziale e \(B(T') \leq B(T)\) per il punto sopra, si ha che \(B(T) = B(T')\),
	ovvero che anche l'albero \(T'\) che soddisfa la proprietà greedy è ottimo.
\end{itemize}

È utile notare che la scelta greedy non distingue se i nodi contenuti nella priority queue siano foglie o radici di
sottoalberi, ma si basa solo sulla loro frequenza. Ciò significa che la dimostrazione sopra (che lavora solo con le foglie)
vale anche per i nodi interni ad un certo livello dell'albero (che eventualmente potrebbero essere sostituiti da foglie
di pari frequenza senza modificare il costo dell'albero).

\newpage

\subsection{(2,1)-boxing}
\subsubsection*{Formalizzazione del problema}
Sia data una sequenza di oggetti \(O = \left\{ o_1, o_2, \dots o_n \right\}\) con la sequenza dei rispettivi pesi
\(P = \left\{ p_1, p_2, \dots p_n \right\}\) positivi e minori di 1,ordinati in peso crescente (\(0 < p_1 \leq p_2 \leq \dots \leq p_n \leq 1\)),
e sia data inoltre una sequenza di scatole \(B = \left\{ b_1, b_2, \dots b_m \right\}\), tali che ogni scatola
\(b_p = \{ o_i, o_j \}\) può contenere al massimo due oggetti \(o_i, o_j\) di peso massimo complessivo inferiore ad 1 (\(p_i + p_j \leq 1\)).

Il problema del (2,1)-boxing consiste nel trovare una disposizione degli oggetti \(O\) nelle scatole \(B\) che minimizzi
il numero di scatole utilizzate.

\subsubsection*{Definizione della scelta greedy}
Si sceglie di inserire per ogni scatola l'oggetto più pesante \(o_j\) non ancora assegnato (l'ultimo della sottostringa
\(O_{i:j}\) degli oggetti non ancora assegnati) con eventuale aggiunta del più leggero \(o_i\) non ancora assegnato (il
primo della sottostringa \(O_{i:j}\)) se il peso lo permette.

Di seguito si formalizza la scelta greedy attraverso la proprietà di sottostruttura espressa in termini di equazioni alle
ricorrenze.
\[B(O_{i:j}, \; P_{i:j}) = \begin{cases}
	\{b_p = \{o_i\}\} & \text{se } i = j \\[3pt]
	\{b_p = \{o_j\}\} \;\; \cup \;\; B(O_{i:j-1}, \; P_{i:j-1}) & \text{se } i < j \text{ e } p_i + p_j > 1 \\[3pt]
	\{b_p = \{o_i, \, o_j\}\} \;\; \cup \;\; B(O_{i+1:j-1}, \; P_{i+1:j-1}) & \text{se } i < j \text{ e } p_i + p_j \leq 1
\end{cases}\]

\subsubsection*{Dimostrazione della correttezza}
Per prima cosa si osserva che la scelta è fattibile, ovvero soddisfa il limite massimo di peso per ogni scatola imposto
dal problema in tutti e tre i casi della definizione ricorsiva della scelta greedy.

Si dimostra ora che esiste una soluzione ottima che soddisfa la scelta greedy.
\begin{itemize}
	\item si ipotizza una soluzione ottima \(B^*(O_{i:j}, \; P_{i:j})\) per l'istanza \(O_{i:j}\), \(P_{i:j}\) e si analizza
	la scatola contenente l'oggetto più pesante \(b_p = \{o_j, \dots\}\):
	\item se la scatola contiene solo l'oggetto più pesante \(b_p = \left\{o_j\right\}\) e vale \(p_i + p_j > 1\), allora
	la scelta greedy è banalmente rispettata
	\item se la scatola contiene solo l'oggetto più pesante \(b_p = \left\{o_j\right\}\), ma vale \(p_i + p_j \leq 1\),
	allora è possibile spostare l'oggetto più leggero \(o_i\) dalla scatola in cui è inserito alla scatola \(b_p\) senza
	violare il vincolo di peso massimo e senza modificare il numero di scatole utilizzate, ottenendo una nuova soluzione
	ottima che rispetta la scelta greedy
	\item se la scatola contiene sia l'oggetto più pesante che quello più leggero \(b_p = \left\{o_i, o_j\right\}\), allora
	la scelta greedy è banalmente rispettata siccome deve necessariamente valere \(p_i + p_j \leq 1\)
	\item se la scatola contiene l'oggetto più pesante e un oggetto che non è il più leggero \(b_p = \left\{ o_j, o_k \right\}\),
	e l'oggetto più leggero si trova in un'altra scatola \(b_q = \{o_i, o_h\}\) allora è possibile scambiare \(o_i\) con \(o_k\)
	senza violare il vincolo di peso massimo (come illustrato di seguito) e senza modificare il numero di scatole utilizzate,
	ottenendo una nuova soluzione ottima that rispetta la scelta greedy
	\[b_p = \{ o_j, o_k\} \;\Rightarrow\; \begin{cases}
		\; p_j + p_k \leq 1 \;\Rightarrow\; p_j + p_i \leq p_j + p_k \leq 1 & \text{per } p_i \leq p_k \text{ con } o_i \text{ il più leggero} \\
		\; p_j + p_k \leq 1 \;\Rightarrow\; p_h + p_k \leq p_j + p_k \leq 1 & \text{per } p_h \leq p_j \text{ con } o_j \text{ il più pesante}
	\end{cases}\]
\end{itemize}


\end{document}
